\documentclass[11pt]{article}

\usepackage[margin=1in]{geometry}
\usepackage{amsmath,amssymb,amsthm,mathtools}
\usepackage{bm}
\usepackage{xcolor}
\usepackage{enumitem}
\usepackage{graphicx}
\usepackage{hyperref}
\usepackage{booktabs}

\hypersetup{
  colorlinks=true,
  linkcolor=blue,
  citecolor=blue,
  urlcolor=blue
}

\newtheorem{theorem}{Theorem}[section]
\newtheorem{proposition}[theorem]{Proposition}
\newtheorem{lemma}[theorem]{Lemma}
\newtheorem{corollary}[theorem]{Corollary}
\newtheorem{remark}[theorem]{Remark}
\numberwithin{equation}{section}

\newcommand{\E}{\mathbb{E}}
\newcommand{\Pbb}{\mathbb{P}}
\newcommand{\R}{\mathbb{R}}
\newcommand{\N}{\mathbb{N}}
\newcommand{\1}{\mathbf 1}
\newcommand{\sech}{\operatorname{sech}}
\newcommand{\cP}{\mathcal P}
\newcommand{\cK}{\mathcal K}
\newcommand{\cR}{\mathcal R}
\newcommand{\dd}{\,\mathrm d}
\newcommand{\norm}[1]{\left\lVert #1\right\rVert}
\newcommand{\ip}[2]{\left\langle #1,#2\right\rangle}
\newcommand{\Tr}{\operatorname{Tr}}
\newcommand{\diag}{\operatorname{diag}}
\newcommand{\sgn}{\operatorname{sgn}}
\newcommand{\Var}{\operatorname{Var}}
\newcommand{\Cov}{\operatorname{Cov}}
\newcommand{\e}{\mathrm e}
\newcommand{\braket}[1]{\left\langle #1\right\rangle}

\title{A quantitative replica-symmetric bound in strict AT region}
\author{
  Seiichiro Kusuoka\thanks{Department of Mathematics and Mathematical Sciences, Graduate School of Science, Kyoto University. E-mail: \texttt{kusuoka@math.kyoto-u.ac.jp}}
  \and
  Shuta Nakajima\thanks{Department of Mathematics, Keio University. E-mail: \texttt{njima@keio.jp}}
}
\date{}

\begin{document}
\maketitle

\begin{abstract}
We consider the Sherrington--Kirkpatrick model with inverse temperature $\beta>0$ and deterministic external
field $h>0$.  Let $q$ be  the replica-symmetric fixed point: 
$$ q=\E\tanh^2\bigl(h+\beta\sqrt q\,Z\bigr),
 \qquad Z\sim N(0,1).$$
We prove that, uniformly on compact subsets of the strict
de Almeida--Thouless region:
$$\beta^2\E\sech^4\bigl(h+\beta\sqrt q\,Z\bigr) <1,$$ 
the overlap satisfies the concentration:
\[
 \E\braket{(R_{12}-q)^2}=O(N^{-1}).
\]
  As a consequence, we
obtain an $O(N^{-1})$ replica-symmetric free-energy correction and
identify the finite-volume replicon susceptibility.  Unlike the approach
of Lopatto \cite{Lopatto}, our proof does not use the identification of
the limiting free energy with the Parisi variational formula.
\end{abstract}

\section{Model and main results}\label{sec:intro}

The Sherrington--Kirkpatrick model is the basic mean-field model of a
disordered Ising system.  In a nonzero external field, its
replica-symmetric order parameter is determined by a scalar fixed-point
equation, while the de Almeida--Thouless condition describes the
stability of that solution against two-replica perturbations.  The
strict inequality in the AT condition is therefore expected to imply
that the overlap has fluctuations of order $N^{-1/2}$ and that the
replica-symmetric free energy is accurate up to order $N^{-1}$. Our purpose is to prove these claims
uniformly on compact subsets of the strict AT region.
\subsection{Model and the strict AT region}

Let $N \in {\mathbb N}$ and $\Sigma_N\coloneqq\{-1,1\}^N$.  For $\beta>0$ and $h>0$, let
\[
 H_N(\sigma)
 \coloneqq
 \frac{\beta}{\sqrt{2N}}\sum_{i,j=1}^Ng_{ij}\sigma_i\sigma_j
 +h\sum_{i=1}^N\sigma_i , \quad \sigma \in \Sigma _N
\]
where $(g_{ij})_{1\le i,j\le N}$ are independent standard normal random variables. 

Let $q(\beta,h)$ be the unique solution of
\begin{equation}
 q=\E\tanh^2\bigl(h+\beta\sqrt q\,Z\bigr),
 \qquad Z\sim N(0,1),
 \label{eq:q}
\end{equation}
and write $q\coloneqq q(\beta,h)$.
Existence and uniqueness for $h>0$ are standard; see, for example,
\cite[Proposition~1.3.8]{Talagrand}.  Continuity follows directly from
uniqueness: if $(\beta_n,h_n)\to(\beta,h)$, every convergent
subsequence of $q(\beta_n,h_n)\in[0,1]$ has, by dominated convergence
in \eqref{eq:q}, a limit solving the fixed-point equation at
$(\beta,h)$.  Uniqueness identifies that limit with $q(\beta,h)$, so
the entire sequence converges.
Set
\begin{equation}
 r\coloneqq\E\tanh^4\bigl(h+\beta\sqrt q\,Z\bigr),
 \label{eq:ra}
\end{equation}
and let
\begin{equation}
 \alpha\coloneqq\beta^2(1-2q+r)
 =\beta^2\E\sech^4\bigl(h+\beta\sqrt q\,Z\bigr).
 \label{eq:alpha}
\end{equation}
We work on a fixed compact set
\begin{equation}
 \cK\Subset\{(\beta,h):\beta>0,\ h>0,\ \alpha(\beta,h)<1\}.
 \label{eq:K}
\end{equation}
Thus, for some $\delta_{\cK}>0$,
\begin{equation}
 1-\alpha\ge\delta_{\cK}
 \qquad\text{on }\cK.
 \label{eq:delta}
\end{equation}
The maps $q,r,\alpha$ are therefore continuous on $\cK$.  In
particular, compactness and $h>0$ imply
\begin{equation}
 0<q_{\cK}\coloneqq\inf_{(\beta,h)\in\cK}q
 \le \sup_{(\beta,h)\in\cK}q<1.
 \label{eq:qcompact}
\end{equation}
Indeed, $q=0$ would contradict \eqref{eq:q} and $h>0$; 
$q<1$ holds because $\tanh^2(h+\beta\sqrt q\,Z)<1$.

\subsection{The replica-symmetric smart path}

For $0\le s\le1$, define the smart-path Hamiltonian
\begin{equation}
 H_{N,s}(\sigma)
 \coloneqq
 \frac{\beta\sqrt s}{\sqrt{2N}}
 \sum_{i,j=1}^Ng_{ij}\sigma_i\sigma_j
 +\sum_{i=1}^N
 \bigl(h+\beta\sqrt{(1-s)q}\,z_i\bigr)\sigma_i,
 \label{eq:path}
\end{equation}
where $(z_i)$ is an independent standard Gaussian family.  We write
$\braket{\cdot}_s$ for the Gibbs expectation, i.e.
\[
\braket{F}_s \coloneqq \frac{1}{Z_{N, s}}\sum _{\sigma \in \Sigma _N} F(\sigma ) e^{H_{N, s}(\sigma)}
\]
for a function $F$ on $\Sigma _N$, where $Z_{N,s} \coloneqq \sum _{\sigma \in \Sigma _N} e^{H_{N,s}(\sigma )}$. For replicas $\sigma^1,\sigma^2,\ldots$ sampled independently from the
Gibbs measure for fixed disorder, put
\[
 R_{ab}\coloneqq\frac1N\sum_{i=1}^N\sigma_i^a\sigma_i^b.
\]
Set
\[
 \nu_s[f]\coloneqq\E\braket{f}_s,\qquad
 Q_{ab}\coloneqq R_{ab}-q.
\]
We also write
\begin{equation}
 \phi_N(s)\coloneqq\frac1N\E\log Z_{N,s}.
 \label{eq:pathfreeenergy}
\end{equation}

\subsection{Main result}

Define
\begin{equation}
 A_s\coloneqq\nu_s[Q_{12}^2],\qquad
 B_s\coloneqq\nu_s[Q_{12}Q_{13}],\qquad
 C_s\coloneqq\nu_s[Q_{12}Q_{34}].
 \label{eq:ABC}
\end{equation}

\begin{theorem}
\label{thm:main}
There is a constant $C_{\cK}<\infty$ such that, for every $N\ge1$,
\[
 \sup_{\substack{(\beta,h)\in\cK\\0\le s\le1}}
 N\,\nu_s[Q_{12}^2]\le C_{\cK}.
\]
Consequently, if
\[
 \phi_N(\beta,h)\coloneqq
 \frac1N\E\log\sum_{\sigma\in\Sigma_N}e^{H_N(\sigma)}
\]
and
\[
 \phi^{\mathrm{RS}}(\beta,h)
 \coloneqq
 \log2+\E\log\cosh(h+\beta\sqrt q\,Z)
 +\frac{\beta^2}{4}(1-q)^2,
\]
then uniformly on $\cK$, 
\[ 
 0\le\phi^{\mathrm{RS}}(\beta,h)-\phi_N(\beta,h)
 \le\frac{C_{\cK}}N.
\]


Moreover, uniformly for $(\beta,h,s)\in\cK\times[0,1]$,
\[
 N\bigl(A_s-2B_s+C_s\bigr)
 =
 \frac{\alpha}{\beta^2(1-s\alpha)}+o_{\cK}(1).
\]
Here and below, $o_{\cK}(1)\to0$ as $N\to\infty$, uniformly over the
 compact set $\mathcal K$.
\end{theorem}
\subsection{A brief history and previous research} 
Spin glasses were introduced to describe magnetic materials in which the interactions between microscopic spins are both disordered and frustrated. Experiments on dilute magnetic alloys in the 1970s showed behavior very different from that of ordinary ferromagnets: a sharp freezing transition, very slow relaxation, and a strong dependence on the history of the sample. These findings suggested that the low-temperature phase could not be described by a single equilibrium state of the usual kind, and they led to probabilistic models with quenched disorder.

One of the first such models was proposed by Edwards and Anderson \cite{EdwardsAnderson1975}. In their model, spins sit on the vertices of a finite-dimensional lattice and interact through random nearest-neighbor couplings. This formulation is physically natural, but its mathematical analysis is very hard. The local interactions compete with each other because of frustration, and the quenched disorder removes the spatial regularity that classical arguments rely on. As a result, even basic questions about the low-temperature phase remain difficult in finite dimensions.

A decisive simplification was proposed by Sherrington and Kirkpatrick \cite{SherringtonKirkpatrick1975}. They replaced the short-range geometry of the Edwards--Anderson model by a mean-field interaction, in which every spin interacts weakly with every other spin. The resulting model, now called the Sherrington--Kirkpatrick (SK) model, keeps both disorder and frustration but is much easier to study by analytic methods. Using the replica method, Sherrington and Kirkpatrick derived a formula for the limiting free energy under the assumption of replica symmetry. The replica method is not rigorous: it starts from identities for integer moments of the partition function and then uses an analytic continuation in the number of replicas. More importantly, the replica-symmetric formula gave a negative entropy at low temperature, which is unphysical. It was soon understood that the replica-symmetric ansatz could not describe the whole phase diagram.

A fundamental step was taken by de Almeida and Thouless \cite{deAlmeidaThouless1978}, who studied the local stability of the replica-symmetric solution. 
Their computation predicted that the replica-symmetric solution is  stable when
\begin{align}
    \label{eq: AT region}
\beta^2\mathbb E\operatorname{sech}^4\bigl(h+\beta\sqrt q\,Z\bigr)\leq 1,
\end{align}
and unstable when the reverse strict inequality holds. This boundary is now called the de Almeida--Thouless line, or AT line. So, the AT condition comes from the physical stability analysis of the replica-symmetric state.

The failure of replica symmetry below the AT line led Parisi to propose a much richer picture of the low-temperature phase \cite{Parisi1979,Parisi1980}. In this picture, the symmetry among replicas is broken in a hierarchical way, and the thermodynamic state is described by an order parameter that gives the distribution of overlaps between typical configurations. This led to the Parisi variational formula for the limiting free energy. The Parisi theory fixed the low-temperature problem of the replica-symmetric solution and gave a single framework that covers both the replica-symmetric and the replica-symmetry-breaking phases.

These ideas have since found use far beyond condensed matter physics. In computer science, random constraint satisfaction problems such as random $K$-SAT and random graph coloring show threshold phenomena similar to replica symmetry breaking, and algorithms based on the cavity method, such as survey propagation, are still among the best known methods for these problems \cite{MezardParisiZecchina2002,MezardMontanari2009}. In information theory and high-dimensional statistics, the free energy of a suitable spin glass model gives the fundamental limits in error-correcting codes, compressed sensing, and inference in graphical models \cite{Nishimori2001,ZdeborovaKrzakala2016}. In machine learning, the Hopfield model of associative memory is a spin glass in which the stored patterns play the role of quenched disorder \cite{Hopfield1982,AmitGutfreundSompolinsky1985}, and spin glass methods have been used to study the loss landscapes and training dynamics of neural networks. More generally, the notions of frustration, energy landscapes, and metastability, first identified in spin glasses, are now standard tools in the study of complex systems, from protein folding to financial networks.

The rigorous mathematical study of the SK model developed along a different route. An early landmark was the work of Aizenman, Lebowitz, and Ruelle \cite{AizenmanLebowitzRuelle1987}, who proved Gaussian fluctuations of the free energy at high temperature and zero external field, using a cluster expansion. At zero external field, replica symmetry is now known to hold for all $\beta<1$, and the critical point $\beta=1$ separates the replica-symmetric and symmetry-breaking regimes. With an external field, finding the full replica-symmetric region turned out to be much harder. Talagrand proved the replica-symmetric formula at sufficiently high temperature with the cavity method, and a simple interpolation argument due to Lata\l{}a gave another proof in the range $\beta \in [0,1/2)$ and $h\in \R$. A recent work of the authors extended Lata\l{}a's argument to a wider range including the region $\beta\leq 1$. These high-temperature results, together with the interpolation methods developed at the same time, form much of the mathematical basis for the quantitative study of the overlap.

In parallel, a different line of work dealt with the limiting free energy at all temperatures. Guerra introduced an interpolation argument that proved the appropriate Parisi bound, and Talagrand then proved the matching opposite bound. Together, these results established the Parisi formula for the SK model. This turned the study of the phase diagram from a problem based on the replica calculus into a rigorous variational problem. In particular, replica symmetry can be characterized through the minimizer of the Parisi functional, and this variational point of view has become one of the main tools for studying the AT transition.

One question remained: does the AT stability condition exactly characterize the replica-symmetric phase of the SK model with a nonzero external field? Several rigorous results proved replica symmetry in regions inside the predicted AT region, or replica symmetry breaking on the other side of the line, but reaching the physical boundary itself remained difficult. Very recently, Lopatto \cite{Lopatto} proved, for a positive uniform external field, that replica symmetry holds throughout the AT region \eqref{eq: AT region} 
 using the variational characterization of the Parisi measure. The de Almeida--Thouless prediction for the replica-symmetric side of the phase diagram is therefore now rigorously established.

The present paper gives a quantitative counterpart of this picture. Instead of asking only whether the thermodynamic limit is replica symmetric, we work under the strict AT condition
 and prove the quantitative estimate stated above. The strict inequality gives a positive stability margin away from the AT boundary, and our argument turns this margin into finite-$N$ control. In this sense, our result complements the characterization of the replica-symmetric phase. 
\subsection{Outline of the proof: from coupled free energy to optimal overlap bound}
\label{subsec:proof-outline}

We briefly describe the proof of the first claim in the main result; the second and third are obtained from this using a standard interpolation argument (see Section~\ref{sec:conclusion}).

For $u \in \mathbb R$, define the constrained
two-replica pressure by
\begin{equation}
 p^{(2)}_{N,s}(u )
  :=
 \frac1{2N}\E\log
 \sum_{\sigma^1,\sigma^2}
 \exp\left(
 H_{N,s}(\sigma^1)+H_{N,s}(\sigma^2)
 +\frac{u  N}{2}Q_{12}^2\right)
\label{eq:constrained-two-replica-pressure}
\end{equation}
The Guerra--Talagrand interpolation gives
\begin{equation}
p_{N,s}^{(2)}(u )
\le
\mathfrak P_s(u )+ o_\cK(1)
\label{eq:GT-constrained-bound-outline}
\end{equation}
uniformly in $(\beta,h)\in\mathcal K$, $s\in[0,1]$, and
$u\in[-1,1]$. Here $\mathfrak P_s(u )$ denotes the
corresponding one-step replica-symmetry-breaking formula. Using Lopatto's argument \cite{Lopatto}, one obtains that as $u \to 0+$,
 \begin{equation}
\mathfrak P_s(u )
= \mathfrak P_s(0)+ o(u ).
\label{eq:global-GT-gap-outline}
\end{equation}

Let \begin{equation*}
 D_N(s)\coloneqq \mathfrak P_s(0)-\phi_N(s).
\end{equation*} 
By Gaussian integration by parts, we have for $u_0>0$ sufficiently small,  
\[
D_N'(s) =  \frac{\beta^2}{4} \nu_s [Q_{12}^2]  \leq \frac{\beta^2}{u_0} ({p_{N,s}^{(2)}(u_0)-\phi_N(s)}) \leq   \frac{\beta^2}{u_0} D_N(s) +  o_\cK(1),
\]
where we have used Jensen's inequality with $p_{N,s}^{(2)}(u_0)-\phi_N(s)= \mathbb E\log \langle e^{u_0 N Q_{12}^2/2} \rangle_s$ for the first inequality, and \eqref{eq:GT-constrained-bound-outline} and \eqref{eq:global-GT-gap-outline} with $o_\cK(1)$ being arbitrarily small for the last inequality.  Since $D_N(0) = 0$, by Gronwall's inequality, we have
$$D_N(s) = o_\cK(1),\quad  p_{N,s}^{(2)}(u_0)-\phi_N(s)= \frac{1}{2N}\mathbb E\log \langle e^{u_0 N Q_{12}^2/2} \rangle_s = o_\cK(1).$$ 
Hence, using a standard Gaussian Lipschitz concentration inequality, one gets the preliminary
overlap concentration estimate
\begin{equation}
\sup_{\substack{(\beta,h)\in\mathcal K\\ 0\le s\le 1}}
\nu_s\left(
|Q_{12}|\ge\delta
\right)
\le
C_{\mathcal K,\delta}
\exp(-c_{\mathcal K,\delta}N),
\qquad \delta>0.
\label{eq:fixed-deviation-overlap-tail-outline}
\end{equation}

To obtain the $N^{-1}$ bound, we use
the cavity bootstrap developed in the high-temperature analysis of
Talagrand \cite[Chapter~1]{Talagrand}: Estimating the influence of adding one spin, we reach
\begin{equation}
\nu_s[Q_{12}^2] \leq C_\cK 
 \left(
\nu_s|Q_{12}|^3+N^{-1}
\right).
\label{eq:cavity-remainder-outline}
\end{equation}
The preliminary concentration estimate \eqref{eq:fixed-deviation-overlap-tail-outline} yields that for any fixed
$0<\delta$,
\begin{align}
\nu_s|Q_{12}|^3
&\le
\delta\nu_s[Q_{12}^2]
+
8\nu_s\left(|Q_{12}|\ge\delta\right)
\notag\\
&\le
\delta \nu_s[Q_{12}^2]+
C_{\mathcal K,\delta}e^{-c_{\mathcal K,\delta}N}.
\label{eq:cubic-bootstrap-outline}
\end{align}
Consequently, we have
\[
\nu_s[Q_{12}^2] \leq \delta C_\cK \nu_s[Q_{12}^2] +  C_\cK N^{-1} + C_\cK C_{\cK,\delta}e^{-c_{\mathcal K,\delta}N}.
\]
Choosing $\delta=1/(2C_\cK)>0$, we obtain
\begin{equation}
\sup_{\substack{(\beta,h)\in\mathcal K\\ 0\le s\le 1}}
\nu_s[Q_{12}^2]
\le
\frac{C'_{\mathcal K}}{N},
\label{eq:optimal-overlap-bound-outline}
\end{equation}
with some $C'_\cK$. This completes the proof of the quantitative bound for overlaps.
\section{Replica-symmetric coercivity}\label{sec:coercivity}

\subsection{Gaussian calculus}\label{subsec:gaussian}

We record Gaussian integration by parts, which is used below.  Let
$G\coloneqq(G_1,\ldots,G_m)$ be centered Gaussian with covariance
matrix $C$.
If $F$ is continuously differentiable with bounded first derivatives,
then Gaussian integration by parts (see \cite[Vol. 1, Sec. A.4]{Talagrand}) gives
\begin{equation}
 \E[G_iF(G)]=\sum_{j=1}^m C_{ij}\E[\partial_jF(G)].
 \label{eq:GIBP}
\end{equation}

We also use the following standard consequence of Gaussian concentration.
Let \(G\) be a standard Gaussian vector, and let $\{ F_v\} _{v\in I}$ be a finite set of $L$-Lipschitz functions on the vector space.
Then, it holds that
\begin{equation}
 \E\max_{v\in I}\{F_v (G) -\E F_v (G)\}
 \le L\sqrt{2\log|I|}.
 \label{eq:gaussianmaxgeneral}
\end{equation}
Indeed, the Gaussian concentration inequality (see, for example,
\cite{Ledoux}) gives, for every \(t>0\)
and every \(v\in I\),
\[
 \E\exp\bigl(t(F_v(G) -\E F_v (G))\bigr)
 \le \exp\left(\frac{t^2L^2}{2}\right).
\]
This inequality and Jensen's inequality imply
\[ \begin{aligned}
\E\max_{v\in I}\bigl(F_v(G)-\E F_v(G)\bigr) &\le \frac1t\log \E\exp\left( t\max_{v\in I}\bigl(F_v(G)-\E F_v(G)\bigr) \right)\\ &\le \frac1t\log \sum_{v\in I} \E\exp\left(t\bigl(F_v(G)-\E F_v(G)\bigr)\right)\\
&\le \frac{\log|I|}{t}+\frac{tL^2}{2}, \end{aligned} \] 
where the first inequality follows from Jensen's inequality. 
For $|I|\ge2$ and $L>0$, optimize the upper bound by choosing
$t=\sqrt{2\log|I|}/L$.  The cases $|I|=1$ and $L=0$ are immediate.
Thus, we obtain \eqref{eq:gaussianmaxgeneral}.

\subsection{The replica-symmetric path}\label{subsec:ATsign}

Fix $(\beta,h,s)\in\cK\times[0,1]$.
For $t\ge0$, let $\mathsf H_t$ be the heat semigroup, i.e. 
\[
 (\mathsf H_t\varphi)(x)\coloneqq\E\varphi(x+\sqrt t\,Z) .
\]
Then, it holds that
\begin{equation}\label{eq:genH}
\partial _t (\mathsf H_t\varphi)(x) = \frac{1}{2} \mathsf H_t (\partial ^2 \varphi)(x) .
\end{equation}
For $0\le u\le1$ and $x\in\mathbb R$, define
\begin{equation}
 \Psi(u,x)
 \coloneqq
 \E\log\cosh\left(
 x+\beta\sqrt{s(q-u)_+}\,Z
 \right)
 +\frac{s\beta^2}{2}\bigl(1-\max\{u,q\}\bigr),
 \qquad Z\sim N(0,1),
 \label{eq:Psi-explicit}
\end{equation}
where $(a)_+\coloneqq\max\{a,0\}$.  Equivalently,
\[
 \Psi(u,x)
 =
 \begin{cases}
  (\mathsf H_{s\beta^2(q-u)}\log\cosh)(x)
  +\dfrac{s\beta^2}{2}(1-q),
  &0\le u\le q,\\[2mm]
  \log\cosh(x)+\dfrac{s\beta^2}{2}(1-u),
  &q\le u\le1.
 \end{cases}
\]
Differentiating \eqref{eq:Psi-explicit} in $x$ gives
\begin{equation}
 \partial_x\Psi(u,x)
 =
 \E\tanh\left(
 x+\beta\sqrt{s(q-u)_+}\,Z
 \right)
 =
 \begin{cases}
  (\mathsf H_{s\beta^2(q-u)}\tanh)(x),
      &0\le u\le q,\\
  \tanh(x),&q\le u\le1.
 \end{cases}
 \label{eq:Psi-x}
\end{equation}

For the upper interval, introduce the tilted heat semigroup
\begin{equation}
 (\mathsf T_t\varphi)(x)
 \coloneqq e^{-t/2}
 \frac{\bigl(\mathsf H_t
   (\varphi \cdot \cosh )\bigr)(x)}{\cosh(x)}.
 \label{eq:tiltedsemigroup}
\end{equation}
Differentiation under the Gaussian integral shows that its generator is
$\frac12\partial_{xx}+\tanh(x)\partial_x$.  For each $u\in[0,1]$,
let $X_{s,u}$ denote a random variable whose law is specified by
\begin{equation}
 \E\varphi(X_{s,u})
 \coloneqq
 \begin{cases}
  \bigl(\mathsf H_{\beta^2(1-s)q+s\beta^2u}\varphi\bigr)(h),
      &0\le u\le q,\\[1mm]
  \bigl(\mathsf H_{\beta^2q}
      (\mathsf T_{s\beta^2(u-q)}\varphi)\bigr)(h),
      &q\le u\le1.
 \end{cases}
 \label{eq:localfieldlaw}
\end{equation}
In particular,
\begin{equation}\label{eq:Xsq}
X_{s,q}\stackrel{\mathrm d}=h+\beta\sqrt q\,Z.
\end{equation}
The replica-symmetric free energy for $H_{N,s}$ along the path is
\begin{align}
 P_s^*& \coloneqq\log2+
 \E\Psi(0,h+\beta\sqrt{(1-s)q}\,Z)
 -\frac{s\beta^2}{2}\int_q^1u\,\dd u
 \label{eq:scalartrial}\\
 &= \log2+\E\log\cosh(h+\beta\sqrt q\,Z)
 +\frac{s\beta^2}{4}(1-q)^2,
 \label{eq:RSpathvalue}
\end{align}
where we have used the semigroup property and
$\beta^2(1-s)q+s\beta^2q=\beta^2q$.

\begin{lemma}
\label{lem:uppercomparison}
If $s\beta^2(1-q)>1$, then
\begin{equation}
 \beta^2\E\sech^4(X_{s,u})\le\alpha,
 \qquad q\le u\le1.
 \label{eq:uppercomparison}
\end{equation}
\end{lemma}



\begin{proof}
Fix $q\le u\le1$ and set
\[
 \sigma^2\coloneqq\beta^2q,
 \qquad
 t\coloneqq s\beta^2(u-q),
 \qquad
 X\coloneqq h+\sigma Z.
\]
By \eqref{eq:localfieldlaw},
\[
 \E\sech^4(X_{s,u})
 =
 \E\bigl[(\mathsf T_t\sech^4)(X)\bigr].
\]
Thus it suffices to prove
\[
 \E\bigl[(\mathsf T_t\sech^4)(X)\bigr]
 \le
 \E\sech^4(X).
\]

We first note that
\begin{equation}
 h<\sigma^2.
 \label{eq:h-less-sigma}
\end{equation}
Indeed, suppose that $h\ge\sigma^2$. Since the convolution of an even
function decreasing on $[0,\infty)$ with a centered Gaussian density has
the same property,
\[
 1-q
 =
 \E\sech^2(h+\sigma Z)
 \le
 \E\sech^2(\sigma^2+\sigma Z).
\]
Let $\varepsilon$ be a symmetric Rademacher variable, independent of $Z$,
and put $Y=\sigma^2\varepsilon+\sigma Z$. Then
$\E[\varepsilon\mid Y]=\tanh Y$ due to $\mathbb P(\varepsilon=\pm 1,\,Y\in dy) = \frac{1}{2\sqrt{2\pi \sigma^2}}e^{-(y - (\pm \sigma^2) )^2/2\sigma^2} dy$, and hence
\[
 \E\sech^2(\sigma^2+\sigma Z)
 =
 \E\operatorname{Var}(\varepsilon\mid Y)
 \le
 \inf_{a\in\R}\E(\varepsilon-aY)^2
 =
 \frac{1}{1+\sigma^2}.
\]
Therefore
\[
 (1-q)(1+\beta^2q)\le1,
\]
which, since $q>0$, implies $\beta^2(1-q)\le1$. This contradicts
$s\beta^2(1-q)>1$ because $s\le1$.

Now put
\[
 f\coloneqq\sech^4,
 \qquad
 g_t\coloneqq\mathsf T_t f,
 \qquad
 M(t)\coloneqq\E g_t(X).
\]
Since
\[
 g_t(x)
 =
 e^{-t/2}\frac{\mathsf H_t(\sech^3)(x)}{\cosh x},
\]
$g_t$ is even and nonincreasing on $[0,\infty)$; in particular,
\[
 g_t'(x)\le0,
 \qquad x>0.
\]
By \eqref{eq:genH} and \eqref{eq:tiltedsemigroup},
\begin{equation}\label{eq:partialT_t}
 \partial_t(\mathsf T_t\varphi)(x)
 =\frac{1}{2\cosh^2x}\partial_x\left(
   \cosh^2x\,\partial_x(\mathsf T_t\varphi)(x)
  \right).
\end{equation}
Let $\phi(x)=(2\pi)^{-1/2}e^{-x^2/2}$ be the standard normal density.
Since $X\sim N(h,\sigma^2)$, its density is
$\sigma^{-1}\phi((x-h)/\sigma)$.
From \eqref{eq:partialT_t} and integration by parts
in divergence form we have
\[
 \begin{aligned}
 M'(t)
 &=\frac12\int_{\mathbb R}
   \partial_x\bigl(\cosh^2x\,g_t'(x)\bigr)
   \frac{\phi((x-h)/\sigma)}{\sigma\cosh^2x}\,\dd x\\
 &=-\frac12\int_{\mathbb R}\cosh^2x\,g_t'(x)
   \partial_x\left(
     \frac{\phi((x-h)/\sigma)}{\sigma\cosh^2x}
   \right)\dd x\\
 &=\E\left[
   \left(\tanh X+\frac{X-h}{2\sigma^2}\right)g_t'(X)
 \right].
 \end{aligned}
\]
Pairing $x$ and $-x$
gives
\begin{equation}\label{eq:M'}
 M'(t)=\int_0^\infty g_t'(x)K(x)\,\dd x,
\end{equation}
where
\[
 \begin{aligned}
 K(x)
 &\coloneqq
 \frac1\sigma\phi\left(\frac{x-h}{\sigma}\right)
 \left(\tanh x+\frac{x-h}{2\sigma^2}\right)+
 \frac1\sigma\phi\left(\frac{-x-h}{\sigma}\right)
 \left(\tanh x+\frac{x+h}{2\sigma^2}\right).
 \end{aligned}
\]
Since
\[
 \frac{\phi((-x-h)/\sigma)}{\phi((x-h)/\sigma)}
 =e^{-2hx/\sigma^2},
\]
we obtain
\[
 \frac{\sigma K(x)}{
  \phi((x-h)/\sigma)+\phi((-x-h)/\sigma)}
 =
 \tanh x+\frac{x}{2\sigma^2}
 -\frac{h}{2\sigma^2}\tanh\left(\frac{hx}{\sigma^2}\right).
\]
Since $0<h/\sigma^2<1$ (see \eqref{eq:h-less-sigma}),
\(
 \tanh\left(\frac{hx}{\sigma^2}\right)\le\tanh x .
\)
This inequality and \eqref{eq:K} implis that for $x>0$
\[
 \frac{\sigma K(x)}{
  \phi((x-h)/\sigma)+\phi((-x-h)/\sigma)}
 \ge
 \left(1-\frac{h}{2\sigma^2}\right)\tanh x
 +\frac{x}{2\sigma^2}
 >0.
\]
Thus $K(x)>0$, while $g_t'(x)\le0$, and therefore \eqref{eq:M'} yields
\[
 M'(t)\le0.
\]
In particular, it holds that $M(t)\le M(0)$. From this inequality we obtain
\[
 \beta^2\E\sech^4(X_{s,u})
 \le
 \beta^2\E\sech^4(h+\beta\sqrt q\,Z)
 =\alpha.
\]
\end{proof}

\begin{proposition}
\label{prop:pathRS}
The function $g_s$ defined by
\begin{equation}
 g_s(u)\coloneqq\E\bigl[(\partial_x\Psi(u,X_{s,u}))^2\bigr]
 \label{eq:gs-def}
\end{equation}
satisfies
\begin{equation}
 g_s(u)-u
 \begin{cases}
 >0,&0\le u<q,\\
 =0,&u=q,\\
 <0,&q<u\le1.
 \end{cases}
 \label{eq:ATsign}
\end{equation}
Moreover, there are constants $c_{\cK},\varepsilon_{\cK}>0$ such that
\begin{equation}
 |g_s(u)-u|\ge c_{\cK}|u-q|
 \qquad\text{whenever }|u-q|\le\varepsilon_{\cK}.
 \label{eq:linearATsign}
\end{equation}
\end{proposition}

\begin{proof}
Since \eqref{eq:Psi-x} and \eqref{eq:Xsq} imply
\[
g_s(q) = \E\bigl[ \tanh ^2 (X_{s,q}) \bigr] =q ,
\]
\eqref{eq:linearATsign} holds for $u=q$.

We consider the case that $u<q$. 
Since \eqref{eq:Psi-x} gives
\[
g_s(u) = \E\left[ \bigl(\mathsf H_{s\beta^2(q-u)} (\tanh )\bigr)(X_{s,u}) ^2 \right] ,
\]
by noting that \eqref{eq:genH} and \eqref{eq:localfieldlaw} imply
\[
\partial _u \E\varphi(X_{s,u}) = \frac{s\beta ^2}{2} \bigl(\mathsf H_{\beta^2(1-s)q+s\beta^2u} (\partial ^2 \varphi ) \bigr) (h) = \frac{s\beta ^2}{2} \E\left[ (\partial ^2 \varphi ) (X_{s,u}) \right] ,
\]
we have
\begin{align*}
g_s'(u) &= -s\beta^2 \left. \partial _v \E\left[ \bigl(\mathsf H_v (\tanh )\bigr)(X_{s,u}) ^2 \right] \right| _{v=s\beta^2(q-u)} + \left. \partial _v \E\left[ \bigl(\mathsf H_{s\beta^2(q-u)} (\tanh )\bigr)(X_{s,v}) ^2 \right] \right| _{v=u} \\
%
&= - s\beta^2 \E\left[ \bigl(\mathsf H_{s\beta^2(q-u)} (\tanh )\bigr)(X_{s,u}) \bigl(\mathsf H_{s\beta^2(q-u)} ( \partial ^2 \tanh )\bigr)(X_{s,u}) \right] \\
&\quad + \frac{s\beta ^2}{2} \E\left[ \left. \left( \partial _x ^2 \bigl(\mathsf H_{s\beta^2(q-u)} (\tanh ) (x)\bigr) ^2 \right) \right| _{x=X_{s,u}} \right] \\
%
&= s\beta ^2 \E\left[ \bigl(\mathsf H_{s\beta^2(q-u)} (\partial \tanh )\bigr) (X_{s,u}) ^2 \right] .
\end{align*}
Thus, we obtain
\[
 g_s'(u)
 =s\beta^2\E\left[
   \left\{
    \bigl(\mathsf H_{s\beta^2(q-u)}
      ( \sech^2 )\bigr)(X_{s,u})
   \right\}^2
 \right].
\]
Jensen's inequality, the semigroup property and \eqref{eq:localfieldlaw} imply
\[
 g_s'(u)
 \le s\beta^2\E\left[
   \bigl(\mathsf H_{s\beta^2(q-u)}
     ( \sech^4 )\bigr)(X_{s,u})
 \right]
 =s\beta^2\E\sech^4(h+\beta\sqrt q\,Z)
 =s\alpha.
\]
Since $g_s(q)=q$, the integration of both sides implies
\[
q - g_s(u) \leq s \alpha (q-u) .
\]
Hence, it holds that
\begin{equation}
 g_s(u)-u\ge(1-s\alpha)(q-u)
 \qquad(0\le u<q).
 \label{eq:leftstrict}
\end{equation}
This yields \eqref{eq:ATsign} for $0\le u<q$.

We consider the case that $u\ge q$.
In this case, \eqref{eq:Psi-x} gives $g_s(u)=\E\tanh^2(X_{s,u})$.
On the other hand, \eqref{eq:genH}, \eqref{eq:tiltedsemigroup} and \eqref{eq:localfieldlaw} imply
\begin{align*}
\partial _u \E\varphi(X_{s,u}) &= \partial _u \mathsf H_{\beta ^2 q} \left( e^{-\frac{s\beta ^2}{2}(u-q)} \frac{\mathsf H_{s\beta ^2 (u-q)}(\varphi \cdot \cosh )}{\cosh} \right) (h) \\
%
&= -\frac{s\beta ^2}{2}\E\varphi(X_{s,u}) + \frac{s\beta ^2}{2} \mathsf H_{\beta ^2 q} \left( e^{-\frac{s\beta ^2}{2}(u-q)} \frac{\mathsf H_{s\beta ^2 (u-q)}( \partial ^2 (\varphi \cdot \cosh) )}{\cosh} \right) (h) \\
%
&= \frac{s\beta ^2}{2} \E (\partial ^2 \varphi )(X_{s,u}) + s\beta ^2 \mathsf H_{\beta ^2 q} \left( e^{-\frac{s\beta ^2}{2}(u-q)} \frac{\mathsf H_{s\beta ^2 (u-q)}( ( \partial \varphi ) \cdot \sinh )}{\cosh} \right) (h).
\end{align*}
Hence,
\begin{align*}
g_s'(u) &= s\beta ^2 \E \sech ^4 (X_{s,u}) -2 s\beta ^2 \E \tanh ^2(X_{s,u}) \sech ^2 (X_{s,u})\\
&\quad + 2 s\beta ^2 \mathsf H_{\beta ^2 q} \left( e^{-\frac{s\beta ^2}{2}(u-q)} \frac{\mathsf H_{s\beta ^2 (u-q)}( \sech \cdot \tanh ^2 )}{\cosh} \right) (h)
\end{align*}
Thus, by recalling \eqref{eq:tiltedsemigroup} and \eqref{eq:localfieldlaw}, we have
\begin{equation}
 g_s'(u)=s\beta^2\E\sech^4(X_{s,u}),
 \qquad g_s'(q)=s\alpha<1.
 \label{eq:gs-upper-derivative}
\end{equation}

If $s\beta^2(1-q)>1$, Lemma~\ref{lem:uppercomparison} gives
$g_s'(u)\le s\alpha <1$, so $g_s(u)<u$ for $u>q$.
We consider the case that $s\beta^2(1-q)\le1$.
Since
$\sech^4(x)\le\sech^2(x)=1-\tanh^2(x)$ for every $x\in\mathbb R$,
by \eqref{eq:gs-upper-derivative} we have
\begin{equation}\label{eq:gderivative}
 g_s'(u)\le s\beta^2(1-g_s(u)).
\end{equation}
Then, \eqref{eq:gderivative} implies that for $u>q$
\[
\partial _u (g_s(u) -u) \leq -s\beta ^2 (g_s(u)-u) + s\beta ^2 (1-u) -1 \leq -s\beta ^2 (g_s(u)-u) .
\]
Hence, by noting that \eqref{eq:gs-upper-derivative} implies $g_s(u)-u <0$ for $u>q$ sufficiently close to $q$ and by applying Gr\"onwall's inequality we have
\[
g_s(u) - u < 0 \qquad (u>q) .
\]
This yields \eqref{eq:ATsign} for $u>q$.

We have \eqref{eq:linearATsign} from \eqref{eq:leftstrict} and \eqref{eq:gs-upper-derivative}.
\end{proof}

\subsection{Two-replica Guerra--Talagrand coercivity}\label{subsec:GT}

For an attainable overlap $v\in\mathcal R_N
\coloneqq\{-1,-1+2/N,\ldots,1\}$, define
\begin{equation}
 Z^{(2)}_{N,s}(v)
 \coloneqq
 \sum_{\substack{\sigma^1,\sigma^2\in\Sigma_N\\R_{12}=v}}
 e^{H_{N,s}(\sigma^1)+H_{N,s}(\sigma^2)}
 \label{eq:constrainedZ}
\end{equation}
and retain the notation $P_s^*$ from \eqref{eq:RSpathvalue}.  For
$\lambda\in\mathbb R$, define
\begin{equation}
 f_\lambda(x_1,x_2)
 \coloneqq\log\left(\frac14\sum_{\varepsilon_1,\varepsilon_2=\pm1}
 \exp(\varepsilon_1x_1+\varepsilon_2x_2
       +\lambda\varepsilon_1\varepsilon_2)\right).
 \label{eq:GTterminal}
\end{equation}

For $(\beta,h,s)\in\cK\times[0,1]$, define $\xi_s$ by
\begin{equation}
 \xi_s(r)\coloneqq\beta^2(1-s)qr+\frac{s\beta^2}{2}r^2.
\label{eq:xi}
\end{equation}
Given $v\in[-1,1]$, write
$r_v\coloneqq|v|$ and put $\iota_v\coloneqq1$ for $v\ge0$ and
$\iota_v\coloneqq-1$ for $v<0$.  Let $\{ a_j \}_{j=1}^{M+1}$ be the increasing list of the distinct elements of $\{0,q,r_v,1\}$, i.e.
\[
\{ a_j \}_{j=1}^{M+1}= \{0,q,r_v,1\} \quad and \quad a_j< a_{j+1} \ (0\le j \le M).
\]
Note that $M=1$ if $r_v =q,1$, and $M=2$ otherwise.
Define
\begin{equation}
 Q_j^v\coloneqq
 \begin{pmatrix}
 a_j&\iota_v(a_j\wedge r_v)\\
 \iota_v(a_j\wedge r_v)&a_j
 \end{pmatrix},
 \qquad 0\le j\le M+1.
\end{equation}
For $1\le j\le M$, set
\begin{equation}
 m_j\coloneqq
 \begin{cases}
 0,&a_j\le q,\\[1mm]
 \frac12,&q<a_j\le r_v,\\[1mm]
 1,&a_j>\max\{q,r_v\}.
 \end{cases}
\end{equation}
The levels $\{ m_j\} _{j=1}^{M}$ are nondecreasing.  More explicitly, if $r_v\le q$, $\{ m_j\} _{j=1}^{M} = \{ 0,1\}$; if $q<r_v<1$, $\{ m_j\} _{j=1}^{M} = \{ 0,\frac12,1\}$; and if $r_v=1$, $\{ m_j\} _{j=1}^{M} = \{ 0,\frac12\}$.  Thus
$\{ m_j\} _{j=1}^{M} \subset [0,1]$ in every case, including
the cases that $r_v=0$ and that $r_v=q$.  In the interpolation below we separately set
$m_{M+1}=1$, even in the case that $r_v=1$.

Define, entrywise,
\begin{equation}
 B_j\coloneqq\beta^2(1-s)q
 \begin{pmatrix}1&1\\1&1\end{pmatrix}
 +s\beta^2Q_j^v
 =\xi_s'(Q_j^v),
 \qquad 0\le j\le M+1.
\end{equation}
Since $r_v$ is a breakpoint, every increment is of one of the forms
\begin{equation}
 Q_j^v-Q_{j-1}^v=
 \begin{cases}
 (a_j-a_{j-1})
 \begin{pmatrix}1&\iota_v\\\iota_v&1\end{pmatrix},&a_j\le r_v,\\[3mm]
 (a_j-a_{j-1})I_2,&a_{j-1}\ge r_v.
 \end{cases}
 \label{eq:GT-overlap-increments}
\end{equation}
Both displayed matrices are positive semidefinite.  Consequently,
\begin{equation}\label{eq:finite-step-vector-increments}
 B_j-B_{j-1}=s\beta^2(Q_j^v-Q_{j-1}^v)\succeq0.
\end{equation}

For a positive semidefinite $2\times2$ matrix $C$, let $G_C$ be a
centered Gaussian vector with covariance $C$, and define
\begin{equation}
 \mathcal T_{m,C}F(x)\coloneqq
 \begin{cases}
 m^{-1}\log\E e^{mF(x+G_C)},&m>0,\\[2mm]
 \E F(x+G_C),&m=0.
 \end{cases}
 \label{eq:GTsemigroup}
\end{equation}
Define deterministic stage functions by
\begin{equation}
 U_{s,M}^{\lambda,v}=f_\lambda,
 \qquad
 U_{s,j-1}^{\lambda,v}
 =\mathcal T_{m_j,B_j-B_{j-1}}U_{s,j}^{\lambda,v},
 \quad ( 1\le j \le M),
 \label{eq:U-stage-recursion}
\end{equation}
and for $x=(x_1,x_2) \in {\mathbb R}^2$ set
\begin{equation}
 U_s^{\lambda,v}(x)\coloneqq U_{s,0}^{\lambda,v}(x).
\end{equation}
When a breakpoint should be clarified, we use the discrete notation
\[
 U_s^{\lambda,v}(a_j,x)\coloneqq U_{s,j}^{\lambda,v}(x).
\]

Equivalently, choose independent centered Gaussian vectors
$G_1,\ldots,G_M$ with $\Var(G_j)=B_j-B_{j-1}$ and let
$\mathcal F_j=\sigma(G_1,\ldots,G_j)$, with $\mathcal F_0$ trivial.
For fixed $x\in\mathbb R^2$, put
\[
 X_M^\lambda=f_\lambda\left(x+\sum_{k=1}^MG_k\right) ,
\]
and for $0\le j \le M$ define $X_j$ recursively by
\begin{equation}
 X_{j-1}^\lambda \coloneqq
 \begin{cases}
 \displaystyle\frac1{m_j}\log\E\left[
 e^{m_jX_j^\lambda}\mid\mathcal F_{j-1}\right],&m_j>0,\\[3mm]
 \displaystyle\E\left[X_j^\lambda\mid\mathcal F_{j-1}\right],&m_j=0.
 \end{cases}
 \label{eq:U-conditional-recursion}
\end{equation}
Then, it holds that
\begin{equation}\label{eq:relationXU}
 X_j^\lambda=U_{s,j}^{\lambda,v}
 \left(x+\sum_{k=1}^jG_k\right),
 \qquad
 X_0^\lambda = U_s^{\lambda,v}(x).
\end{equation}
In particular, $U_s^{\lambda,v}(x)$ is deterministic.  This definition
also covers the case that $s=0$.
In the case, every covariance increment vanishes.

Next, put
\begin{equation}
 d_j\coloneqq\frac{s\beta^2}{2}
 \sum_{k,\ell =1}^2(Q_{j,k\ell}^v)^2,
 \qquad 0\le j\le M
\end{equation}
where $Q_{j,k\ell}^v$ is the $(k,\ell)$ entry of $Q_{j}^v$.
As a function of a breakpoint value $a$, this quantity is given by
\[
 d(a)=
 \begin{cases}
 2s\beta^2a^2,&0\le a\le r_v,\\[1mm]
 s\beta^2(a^2+r_v^2),&r_v\le a\le1.
 \end{cases}
\]
Indeed, it holds that $d_j = d(a_j)$.
Since $d(a)$ is continuous and nondecreasing, 
\begin{equation}\label{eq:finite-step-scalar-increments}
d_j-d_{j-1}\ge0.
\end{equation}
Define
\begin{equation}\label{eq:defLv}
 \mathcal L_s(v)\coloneqq
 \frac12\sum_{j=1}^Mm_j(d_j-d_{j-1}).
\end{equation}
This normalization is the scalar Gaussian contribution at the
$t=1$ endpoint of the interpolation in Appendix~\ref{app:specialGT}, and is independent of $v$.  Indeed, if $r_v\le q$, then
\[
\sum_{j=1}^Mm_j(d_j-d_{j-1}) = d(1)-d(q)=s\beta^2(1-q^2).
\]
If $q<r_v < 1$,
\begin{align*}
&\sum_{j=1}^Mm_j(d_j-d_{j-1})
 =\frac12(d(r_v)-d(q))+d(1)-d(r_v) \\
& =s\beta^2(r_v^2-q^2)+s\beta^2(1-r_v^2)
 =s\beta^2(1-q^2).
\end{align*}
If $r_v=1$,
\[
 \sum_{j=1}^M m_j(d_j-d_{j-1}) = \frac12(d(1)-d(q))=s\beta^2(1-q^2).
\]
Therefore,
\begin{equation}
 \mathcal L_s(v)=\frac{s\beta^2}{2}(1-q^2),
 \label{eq:GTcorrection}
\end{equation}
in particular $\mathcal L_s(v)$ is independent of $v$.

With $Y_0=h+\beta\sqrt{(1-s)q}\,Z_0$, set
\begin{equation}
 \mathfrak P_s(\lambda,v)
 \coloneqq2\log2+\E U_s^{\lambda,v}(Y_0,Y_0)
 -\lambda v-\mathcal L_s(v).
 \label{eq:GTfunctional}
\end{equation}

The following finite-volume Guerra--Talagrand upper bound is
Lemma~\ref{lem:specialGT}, proved in Appendix~\ref{app:specialGT}. It is a special case
of the general two-dimensional bound of Talagrand and Chen
\cite{Talagrand,ChenGT}:
\begin{equation}
 \frac1N\E\log Z^{(2)}_{N,s}(v)
 \le \inf_{\lambda\in \mathbb R} \mathfrak P_s(\lambda,v).
 \label{eq:GTfunctionalbound}
\end{equation}

The following technical lemma is also postponed to Appendix~\ref{app: technical}.

\begin{lemma}
\label{lem:GTcontinuity}
On every compact set on which $\lambda$ is bounded, the maps
\[
 (\beta,h,s,\lambda,v)\longmapsto
 \mathfrak P_s(\lambda,v)-2P_s^*,
 \qquad
 (\beta,h,s,\lambda,v)\longmapsto
 \partial_\lambda\mathfrak P_s(\lambda,v)
\]
are jointly continuous for
$(\beta,h,s,v)\in\cK\times[0,1]\times[-1,1]$.
\end{lemma}

\begin{lemma}
\label{lem:U-finite-step-derivatives}
In the recursion \eqref{eq:U-conditional-recursion}, define the tilted
conditional expectation by
\[
\E_j^\lambda[H]
\coloneqq
\frac{
\E\left[
H e^{m_jX_j^\lambda}\mid\mathcal F_{j-1}
\right]
}{
\E\left[
e^{m_jX_j^\lambda}\mid\mathcal F_{j-1}
\right]
},
\qquad m_j>0,
\]
and let $\E_j^\lambda[H]=\E[H\mid\mathcal F_{j-1}]$ when $m_j=0$.
Then
\begin{align}
\partial_\lambda X_{j-1}^\lambda
&=
\E_j^\lambda
\left[
\partial_\lambda X_j^\lambda
\right],
\label{eq:U-first-derivative-recursion}\\
\partial_{\lambda\lambda}X_{j-1}^\lambda
&=
\E_j^\lambda
\left[
\partial_{\lambda\lambda}X_j^\lambda
\right]
+
m_j
\operatorname{Var}_j^\lambda
\left(
\partial_\lambda X_j^\lambda
\right).
\label{eq:U-second-derivative-recursion}
\end{align}
\end{lemma}


\begin{proof}
For $m_j>0$, recalling that $$ X_{j-1}^\lambda= \displaystyle\frac1{m_j}\log\E\left[
 e^{m_jX_j^\lambda}\mid\mathcal F_{j-1}\right],$$
 we have 
\begin{align*}
 \partial_\lambda X_{j-1}^\lambda
 &=\frac{
\E\left[
\partial_\lambda X_j^\lambda e^{m_jX_j^\lambda}\mid\mathcal F_{j-1}
\right]
}{
\E\left[
e^{m_jX_j^\lambda}\mid\mathcal F_{j-1}
\right]}
= \E_j^\lambda[\partial_\lambda X_j^\lambda],\\
 \partial_{\lambda\lambda}X_{j-1}^\lambda
 &=\E_j^\lambda[\partial_{\lambda\lambda}X_j^\lambda]
   +m_j\operatorname{Var}_j^\lambda
       (\partial_\lambda X_j^\lambda).
\end{align*}
For $m_j=0$, direct differentiation of the ordinary conditional
expectation gives the same identities with no variance term.  They complete the proof. 
\end{proof}

\begin{corollary}\label{cor: technical lemma for U and Psi}
It holds
    \begin{equation}
\left|\partial_\lambda U_s^{\lambda,v}(x)\right|\le1,
\qquad
0\le
\partial_{\lambda\lambda}U_s^{\lambda,v}(x)
\le\frac52.
\label{eq:U-lambda-derivative-bounds}
\end{equation}
Consequently, uniformly in $(\beta,h,s,v,\lambda)$,
\begin{equation}
0\leq \partial_{\lambda\lambda}\mathfrak P_s(\lambda,v)\leq \frac52.
 \label{eq:GTsecondderivative-simple}
\end{equation}
\end{corollary}
\begin{proof}
Let $\langle\cdot\rangle_{\lambda,x}$ denote expectation on
$\{\pm1\}^2$ with weights proportional to
\[
 \exp(\varepsilon_1x_1+\varepsilon_2x_2
             +\lambda\varepsilon_1\varepsilon_2) \quad ( (\varepsilon _1 , \varepsilon _2) \in \{ \pm 1 \} ^2 ).
\]
Since $(\varepsilon_1\varepsilon_2)^2=1$,
\[
 \partial_\lambda f_\lambda(x)
 =\langle\varepsilon_1\varepsilon_2\rangle_{\lambda,x},
 \qquad
 \partial_{\lambda\lambda}f_\lambda(x)
 =1-\langle\varepsilon_1\varepsilon_2\rangle_{\lambda,x}^2.
\]
Hence, we have $|\partial_\lambda f_\lambda|\le1$ and
$0\le\partial_{\lambda\lambda}f_\lambda\le1$.
By applying \eqref{eq:U-first-derivative-recursion} recursively, we have
\[
 \operatorname{Var}_j^\lambda(\partial_\lambda X_j^\lambda)
 \le \E_j^\lambda[(\partial_\lambda X_j^\lambda)^2]
 \le1.
\]
Similarly, applying \eqref{eq:U-second-derivative-recursion} recursively, we have 
\[
\partial_{\lambda\lambda}X_0^\lambda = \E_0^\lambda \cdots \E_M^\lambda
\left[
\partial_{\lambda\lambda}X_M^\lambda
\right]
+
\sum _{j=0}^M m_j
\E_0^\lambda \cdots \E_j^\lambda \left[ \operatorname{Var}_j^\lambda
\left(
\partial_\lambda X_j^\lambda
\right) \right] .
\]
This equality, \eqref{eq:relationXU} and the fact that $0\le\partial_{\lambda\lambda}f_\lambda\le1$ yield
\begin{equation}\label{eq:estd2U}
 0\le\partial_{\lambda\lambda}U_s^{\lambda,v}(x)
 \le1+\sum_{j=1}^Mm_j.
\end{equation}
If $r_v\le q$, the sum is at most $1$.  If
$q<r_v<1$, $\{ m_j\} _{j=1}^m = \{ 0,\frac12,1\}$ and their sum is $3/2$.
If $r_v=1$, $\{ m_j\} _{j=1}^m = \{ 0,\frac12\}$ and their sum is $1/2$.
Noting that $M=1,2$, consequently we have in every boundary configuration,
\[
 1+\sum_{j=1}^Mm_j\le\frac52.
\]
Together with \eqref{eq:estd2U}, this proves
\eqref{eq:U-lambda-derivative-bounds}.

Moreover, differentiation under the expectation in
\eqref{eq:GTfunctional}
gives
\[
 \partial_{\lambda\lambda}\mathfrak P_s(\lambda,v)
 =\E\partial_{\lambda\lambda}
   U_s^{\lambda,v}(Y_0,Y_0).
\]
Hence \eqref{eq:U-lambda-derivative-bounds} proves
\eqref{eq:GTsecondderivative-simple}.
\end{proof}
\begin{proposition}
\label{prop:GTcoercivity}
Put
\[
 \Gamma_s(v)\coloneqq\inf_{\lambda\in\mathbb R}
 \mathfrak P_s(\lambda,v).
\]
There is a constant $c_{\cK}>0$ such that
\begin{equation}
 \Gamma_s(v)
 \le 2P_s^*-c_{\cK}(v-q)^2
 \label{eq:Gamma-coercivity}
\end{equation}
for $(\beta,h,s,v)\in\cK\times[0,1]\times[-1,1]$.  Consequently,
for every $N$ and $v\in\mathcal R_N$,
\begin{equation}
 \frac1N\E\log Z^{(2)}_{N,s}(v)
 \le 2P_s^*-c_{\cK}(v-q)^2.
 \label{eq:GTcoercivity}
\end{equation}
Moreover,
\begin{equation}
 \Gamma_s(q)=2P_s^*.
 \label{eq:Gamma-at-q}
\end{equation}
\end{proposition}

\begin{proof}

We first evaluate the finite recursion at $\lambda=0$.  Suppose that
$-q\le v\le q$.  The level-$1$ recursion above $q$ {\color{red}(What is this?)} and the identities
\[
 f_0(x_1,x_2)=\log\cosh x_1+\log\cosh x_2,
 \qquad
 \left.\partial_\lambda f_\lambda(x_1,x_2)
 \right|_{\lambda=0}=\tanh x_1\tanh x_2
\]
give, at the breakpoint $q$, 
\begin{align*}
 U_s^{0,v}(q,x_1,x_2)
 &=s\beta^2(1-q)+\log\cosh x_1+\log\cosh x_2,\\
 \left.\partial_\lambda U_s^{\lambda,v}(q,x_1,x_2)
 \right|_{\lambda=0}
 &=\tanh x_1\tanh x_2.
\end{align*}
All recursions below $q$ have level zero.{\color{red}(What does it mean?)}
Let $(Y_1(v),Y_2(v))$  a Gaussian
pair such that both $Y_1(v)$ and $Y_2(v)$ have means $h$ and variances $\beta^2q$, and have the covariance
\[
 \Cov(Y_1(v),Y_2(v))
 =\beta^2\bigl((1-s)q+sv\bigr).
\]
It follows directly from the conditional expectations below $q$ {\color{red}(What does this mean?)}, \eqref{eq:GTcorrection} and
\eqref{eq:RSpathvalue} that
\begin{equation}
 \mathfrak P_s(0,v)=2P_s^*,
 \qquad
 \partial_\lambda\mathfrak P_s(0,v)= \tilde{g}_s(v)-v \qquad (-q\le v\le q)
 \label{eq:GT-zero-signed-simple}
\end{equation}
where
\[
 \tilde{g}_s(v)
 \coloneqq
 \E[\tanh Y_1(v)\tanh Y_2(v)].
\]
Gaussian differentiation {\color{red}(What is this?)} and the Cauchy--Schwarz inequality yield
\[
 \tilde{g}_s'(v)
 =s\beta^2
  \E[\sech^2Y_1(v)\sech^2Y_2(v)]
 \le s\alpha
\]
where $\alpha$ is defined in \eqref{eq:alpha}.
Since $\Cov(Y_1(q),Y_2(q)) = \beta ^2 q = \operatorname{Var} (Y_0)$, it holds that $\tilde{g}_s (q) = \E[\tanh ^2 (Y_0)] =q$ where $Y_0 = h + \beta \sqrt{(1-s)q}Z_0$.
Hence, it follows that
\begin{equation}
\tilde{g}_s(v) -v
 \ge(1-s\alpha)(q-v)>0,
 \qquad -q\le v<q.
 \label{eq:GT-signed-positive-simple}
\end{equation}

For $0\le v\le q$, put $t=s\beta^2(q-v)$, let
$X\sim N(h,\beta^2((1-s)q+sv))$, and let $Z_1$ and $Z_2$ be random variables with the standard normal law which are also independent of other random variables.  
The variables
\[
 X+\sqrt t\,Z_1,
 \qquad X+\sqrt t\,Z_2,
\]
have the joint law of $(Y_1(v),Y_2(v))$, and are independent under the probability measure conditioned on $X$.
Consequently,
\begin{equation}\label{eq:gv(0<v<q)}
 \tilde{g}_s(v) =\E\left[
 \bigl(\mathsf H_t(\tanh)(X)\bigr)^2
 \right]=g_s(v), \qquad 0\le v \le q
\end{equation}
by \eqref{eq:Psi-x}, \eqref{eq:localfieldlaw}, and \eqref{eq:gs-def}.

Now let $q<v\le1$ and put $t=s\beta^2(v-q)$.  Above $v$ the level is
$1$ and the coordinate increments are independent.  The stage
recursion gives
\begin{align*}
 U_s^{0,v}(v,x_1,x_2)
 &=s\beta^2(1-v)+\log\cosh x_1+\log\cosh x_2,\\
 D_v(x_1,x_2)
 &\coloneqq
 \left.\partial_\lambda U_s^{\lambda,v}(v,x_1,x_2)
 \right|_{\lambda=0}
 =\tanh x_1\tanh x_2.
\end{align*}
For $v=1$, the stage above $v$ is absent and these formulas reduce to
the terminal identities.  Between $q$ and $v$, the increment is
$(\sqrt t\,Z,\sqrt t\,Z)$ and its level is $1/2$.  Hence, on the
diagonal,
\begin{align}
 U_s^{0,v}(q,(x,x))
 &=s\beta^2(1-v)
 +2\log\E\cosh(x+\sqrt t\,Z)\notag\\
 &=2\log\cosh x+s\beta^2(1-q),
 \label{eq:GT-positive-diagonal-value}
\end{align}
where $\E\cosh(x+\sqrt t\,Z)=e^{t/2}\cosh x$.  The derivative
recursion in Lemma~\ref{lem:U-finite-step-derivatives} gives
\begin{align}
 D_q(x,x)
 &=
 \left.\partial_\lambda U_s^{\lambda,v}(q,(x,x))
 \right|_{\lambda=0}\notag\\
 &=\frac{\E[\tanh^2(x+\sqrt t\,Z)
                \cosh(x+\sqrt t\,Z)]}
         {\E\cosh(x+\sqrt t\,Z)}
 =\bigl(\mathsf T_t\tanh^2\bigr)(x).
 \label{eq:GT-positive-diagonal-derivative}
\end{align}
Below $q$, the increment is common to both coordinates and its level
is zero.  Combining it with $Y_0$ in
\eqref{eq:GT-positive-diagonal-value} gives
\[
 \E U_s^{0,v}(Y_0,Y_0)
 =2\E\log\cosh(h+\beta\sqrt q\,Z)+s\beta^2(1-q).
\]
Likewise, \eqref{eq:GT-positive-diagonal-derivative} and
\eqref{eq:localfieldlaw} give
\[
 \E\left.
 \partial_\lambda U_s^{\lambda,v}(Y_0,Y_0)
 \right|_{\lambda=0}
 =\bigl(\mathsf H_{\beta^2q}
       \mathsf T_{s\beta^2(v-q)}\tanh^2\bigr)(h)
 =g_s(v) .
\]
Thus, together with \eqref{eq:GT-zero-signed-simple} and \eqref{eq:gv(0<v<q)}, we obtain
\begin{equation}
 \mathfrak P_s(0,v)=2P_s^*,
 \qquad
 \partial_\lambda\mathfrak P_s(0,v)=g_s(v)-v,
 \qquad 0\le v\le1.
 \label{eq:GT-zero-positive-simple}
\end{equation}

We also need a strict gap in the case that $v<-q$.  Suppose first that $s>0$ and
write
\[
 r=|v|>q,
 \qquad
 t=s\beta^2(r-q)>0.
\]
The stage recursion above $r$, whose level is $1$, gives
\[
 U_s^{0,v}(r,x_1,x_2)
 =s\beta^2(1-r)+\log\cosh x_1+\log\cosh x_2.
\]
The stage recursion from $r$ to $q$, whose level is $1/2$ and whose
Gaussian increment is $(\sqrt t\,Z,-\sqrt t\,Z)$, gives
\[
 U_s^{0,v}(q,x_1,x_2)
 =
 s\beta^2(1-r)
 +2\log\E\sqrt{
   \cosh(x_1+\sqrt t\,Z)\cosh(x_2-\sqrt t\,Z)}.
\]
By Cauchy--Schwarz,
\begin{align*}
 U_s^{0,v}(q,x_1,x_2)
 &\le
 s\beta^2(1-r)+t
 +\log\cosh x_1+\log\cosh x_2\\
 &=\Psi(q,x_1)+\Psi(q,x_2).
\end{align*}
The formulas remain valid at $v=-1$, where the stage above $r$ has
zero covariance.  Since $t>0$, equality holds only when
$x_1=-x_2$.  Below $q$ the recursion is
linear, and it evaluates the two sides at
\[
 X_1=Y_0+\beta\sqrt{sq}\,Z,
 \qquad
 X_2=Y_0-\beta\sqrt{sq}\,Z.
\]
Since $\Pbb(X_1=-X_2)=\Pbb(Y_0=0)=0$, also when $s=1$ because then
$Y_0=h>0$, the propagated inequality is strict.  Therefore
\begin{equation}
 \mathfrak P_s(0,v)<2P_s^*,
 \qquad v<-q,\quad s>0.
 \label{eq:GT-negative-strict-simple}
\end{equation}
At $s=0$, a direct calculation gives
\begin{equation}
 \mathfrak P_0(0,v)=2P_0^*,
 \quad
 \partial_\lambda\mathfrak P_0(0,v)=q-v, \qquad v<-q.
 \label{eq:GT-szero-simple}
\end{equation}

We now prove the desired bound by applying the results above.
Choose
$\varepsilon_{\cK}<q_{\cK}/2$ so that
\eqref{eq:linearATsign} holds when
$|v-q|\le\varepsilon_{\cK}$.  Proposition~\ref{prop:pathRS} gives
$a_{\cK}>0$ such that
\begin{equation}
 (g_s(v)-v)(v-q)
 \le-a_{\cK}(v-q)^2,
 \qquad |v-q|\le\varepsilon_{\cK}.
 \label{eq:GT-local-slope-simple} {\color{red}\mbox{(Why this one is needed?)}}
\end{equation}
These values of $v$ are positive, so
\eqref{eq:GTsecondderivative-simple} and \eqref{eq:GT-zero-positive-simple} imply
\[
 \mathfrak P_s(\lambda,v)
 \le
 2P_s^*+\lambda(g_s(v)-v)
 +\frac{M_{\cK}}2\lambda^2.
\]
Taking $\lambda=a_{\cK}(v-q)/M_{\cK}$ yields
\begin{equation}
 \Gamma_s(v)
 \le
 2P_s^*
 -\frac{a_{\cK}^2}{2M_{\cK}}(v-q)^2,
 \qquad |v-q|\le\varepsilon_{\cK}.
 \label{eq:GT-local-coercivity-simple}
\end{equation}

It remains to consider overlaps bounded away from $q$.  We claim that
\begin{equation}
 \Gamma_s(v)<2P_s^*
 \qquad\text{for }v\ne q.
 \label{eq:GT-pointwise-gap-simple}
\end{equation}
In the case that $0\le v<q$ and $q<v\le1$, this follows from \eqref{eq:GT-zero-positive-simple} and \eqref{eq:ATsign} by moving
$\lambda$ slightly in the direction opposite to the derivative.
For the case $-q\le v<0$, use \eqref{eq:GT-zero-signed-simple} and \eqref{eq:GT-signed-positive-simple}.
For the case that $v<-q$ and $s>0$, apply \eqref{eq:GT-negative-strict-simple}.
For the case that $v<-q$ and $s=0$, \eqref{eq:GT-szero-simple} gives the claim by taking a sufficiently small negative $\lambda$.

Consider the compact set
\[
 \mathcal D_{\cK}
 =
 \left\{
 (\beta,h,s,v):
 (\beta,h)\in\cK,\ 0\le s\le1,\ -1\le v\le1,\
 |v-q|\ge\varepsilon_{\cK}
 \right\}.
\]
At every point $\theta\in\mathcal D_{\cK}$,
\eqref{eq:GT-pointwise-gap-simple} provides
$\lambda_\theta\in\mathbb R$ such that
\[
 \mathfrak P_s(\lambda_\theta,v)<2P_s^*.
\]
In view of Lemma~\ref{lem:GTcontinuity}, this inequality persists with a
positive margin on a neighborhood of $\theta$.
The compactness of ${\mathcal D}_\cK$ gives $\delta _{\cK}>0$ such that
\[
 \Gamma_s(v)\le2P_s^*-\delta _{\cK},
 \qquad |v-q|\ge\varepsilon_{\cK}.
\]
Since $|v-q|\le2$,
\begin{equation}
 \Gamma_s(v)
 \le
 2P_s^*-\frac{\delta _{\cK}}4(v-q)^2,
 \qquad |v-q|\ge\varepsilon_{\cK}.
 \label{eq:GT-far-coercivity-simple}
\end{equation}
Combining \eqref{eq:GT-local-coercivity-simple} and
\eqref{eq:GT-far-coercivity-simple} proves
\eqref{eq:Gamma-coercivity}, with
\[
 c_{\cK}
 =
 \min\left\{
 \frac{a_{\cK}^2}{2M_{\cK}},
 \frac{\delta _{\cK}}4
 \right\}>0.
\]
The finite-volume estimate \eqref{eq:GTcoercivity} follows from
\eqref{eq:GTfunctionalbound}.

Finally, at $v=q$, \eqref{eq:GT-zero-positive-simple} gives
\[
 \mathfrak P_s(0,q)=2P_s^*,
 \qquad
 \partial_\lambda\mathfrak P_s(0,q)=g_s(q)-q=0.
\]
Since \eqref{eq:GTsecondderivative-simple} implies that $\lambda\mapsto\mathfrak P_s(\lambda,q)$ is convex, $\lambda=0$ is
a global minimizer.  Hence
\[
 \Gamma_s(q)=2P_s^*,
\]
which proves \eqref{eq:Gamma-at-q}.
\end{proof}


\section{Preliminary overlap concentration}\label{sec:concentration}

\subsection{Finite-volume Gronwall closure}\label{subsec:gronwall}

For $\rho\ge0$, define the coupled pressure $p^{(2)}_{N,s}(\rho)$ and its normalized excess $F_{N,s}(\rho)$ by
\begin{align}
 p^{(2)}_{N,s}(\rho)
 &\coloneqq
 \frac1{2N}\E\log
 \sum_{\sigma^1,\sigma^2}
 \exp\left(
 H_{N,s}(\sigma^1)+H_{N,s}(\sigma^2)
 +\frac{\rho N}{2}Q_{12}^2\right),
 \label{eq:quadraticpressure}\\
 F_{N,s}(\rho)
 &\coloneqq
 p^{(2)}_{N,s}(\rho)-\phi_N(s)
 =
 \frac1{2N}\E\log
 \braket{e^{\rho NQ_{12}^2/2}}_s ,
 \label{eq:normalizedcoupling}
\end{align}
respectively.
Fix
\begin{equation}
 0<\rho_0<2c_{\cK},
 \label{eq:rho0-choice}
\end{equation}
where $c_{\cK}$ is the coercivity constant in
Proposition~\ref{prop:GTcoercivity}.

\begin{lemma}
\label{lem:sublinearpressure}
It holds
\begin{equation}
 p^{(2)}_{N,s}(\rho_0)\le P_s^*+O_{\cK}(1) 
 \sqrt{\frac{\log(N+1)}N}.
 \label{eq:pressureenvelope}
\end{equation}

\end{lemma}

\begin{proof}
Let $L_v\coloneqq\log Z^{(2)}_{N,s}(v)$.  Each $L_v$ is a Lipschitz function
of the Gaussian disorder with squared Lipschitz constant at most
$O_{\cK}(1)N$.  Indeed,
\[
 \left|\partial_{g_{ij}}L_v\right|
 \le\frac{\sqrt2\beta\sqrt s}{\sqrt N},\qquad
 \left|\partial_{z_i}L_v\right|
 \le2\beta\sqrt{(1-s)q},
\]
and summing up the squares over $i,j$ and $i$ gives the claim.  Hence, by \eqref{eq:gaussianmaxgeneral}, 
\begin{equation}
 \E\max_{v\in\mathcal R_N}(L_v-\E L_v)
 \le O_{\cK}(1)\sqrt{N\log(N+1)}
 \label{eq:gaussianmax}
\end{equation}
where we recall $\mathcal R_N \coloneqq \{ -1,-1+2/N,\ldots,1\}$.
Thus, by using \eqref{eq:gaussianmax},
\eqref{eq:GTcoercivity} and \eqref{eq:rho0-choice}, we have
\begin{align*}
 2Np^{(2)}_{N,s}(\rho_0)
 &=
 \E\log\sum_{v\in\mathcal R_N}
 \exp\left(\E L_v + (L_v - \E L_v ) +\frac{\rho_0N}{2}(v-q)^2\right)\\
 &\le
 \log(N+1)+
 \max_v\left\{\E L_v+\frac{\rho_0N}{2}(v-q)^2\right\}
 +O_{\cK}(1)\sqrt{N\log(N+1)}\\
 &\le2NP_s^*+O_{\cK}(1) 
 \sqrt{{N\log(N+1)}}.
\end{align*}
This proves \eqref{eq:pressureenvelope}.
\end{proof}

Put
\begin{equation}
 D_N(s)\coloneqq P_s^*-\phi_N(s).
 \label{eq:DN}
\end{equation}
\begin{lemma}
\label{lem:preconcentration}
It holds that for $s\in [0,1]$
\begin{equation}
 D_N(s)\le O_{\cK}(1) 
 \sqrt{\frac{\log(N+1)}N},
 \qquad
 A_s\le O_{\cK}(1) 
 \sqrt{\frac{\log(N+1)}N}.
 \label{eq:preconcentration}
\end{equation}
\end{lemma}
\begin{proof}
Write 
$$\epsilon_N:= O_{\cK}(1) 
 \sqrt{\frac{\log(N+1)}N}.$$
Similarly to \cite[(14)]{KusuokaNakajima1}, via Gaussian integration by parts (see \cite[Vol.~1,~Sec.~A.4]{Talagrand}), it follows that
\begin{equation}\label{eq:p-derivative}
\phi_N'(s)
=\frac{\beta^2}{4}\left(
(1-q)^2
-
\nu_s[Q_{12}^2]\right).
\end{equation}
%For $0<s<1$, differentiating under the
%expectation and applying \eqref{eq:GIBP} to the variables %$g_{ij}$ and
%$z_i$ gives
%\begin{align*}
% \phi_N'(s)
% ={}&\frac{\beta^2}{4N^2}\sum_{i,j}
% \E\left[1-\braket{\sigma_i\sigma_j}_s^2\right]
% -\frac{\beta^2q}{2N}\sum_i
% \E\left[1-\braket{\sigma_i}_s^2\right]\\
% ={}&\frac{\beta^2}{4}
% \left((1-q)^2-\nu_s[Q_{12}^2]\right).
%\end{align*}
%Here we used
%\[
% \sum_{i,j}\sigma_i^1\sigma_j^1\sigma_i^2\sigma_j^2
% =N^2R_{12}^2,
% \qquad
% \frac1N\sum_i\E\braket{\sigma_i^1\sigma_i^2}_s
% =\nu_s[R_{12}].
%\]
Combining this identity with \eqref{eq:RSpathvalue} and the definition
of $A_s$ in \eqref{eq:ABC} yields
\begin{equation}
 D_N'(s)=\frac{\beta^2}{4}A_s. 
 \label{eq:DNprime}
\end{equation}
%Both sides extend continuously to $s=0,1$ by finite-dimensional
%Gaussian dominated convergence, so the identity also holds for the
%one-sided endpoint derivatives.
On the other hand, convexity of $F_{N,s}(\rho)$ in $\rho \in [0,\infty)$, where $F_{N,s}(\rho)$ is defined in \eqref{eq:normalizedcoupling}, gives
\begin{equation}\label{eq:lemreconcentration1}
   \frac14A_s
 =\left.\partial_\rho F_{N,s}(\rho)\right|_{\rho=0+}
 \le\frac{F_{N,s}(\rho_0)}{\rho_0}.  
\end{equation}
By \eqref{eq:pressureenvelope},
\begin{equation}\label{eq:lemreconcentration2}
 F_{N,s}(\rho_0)
 \le D_N(s)+\epsilon_N.
\end{equation}
Consequently,
\begin{equation}
 D_N'(s)
 \le\frac{\beta^2}{\rho_0}
 \bigl(D_N(s)+\epsilon_N\bigr).
 \label{eq:gronwallDN}
\end{equation}
Gronwall's inequality, the compactness of $\cK$, the continuity of $D_N (s)$ in $s\in [0,1]$, and the fact that $D_N(0)=0$ imply the inequality for $D_N(s)$ in \eqref{eq:preconcentration}.
The inequality for $A_s$ in \eqref{eq:preconcentration} follows from this inequality, \eqref{eq:lemreconcentration1} and \eqref{eq:lemreconcentration2}.
\end{proof}


\subsection{Fixed-deviation overlap concentration}\label{subsec:tails}

\begin{corollary}
\label{cor:tail}
For every $\varepsilon>0$, there are constants
$c_{\cK,\varepsilon},C_{\cK,\varepsilon}>0$ such that
\begin{equation}
 \sup_{\substack{(\beta,h)\in\cK\\0\le s\le1}}
 \nu_s\bigl[|Q_{12}|\ge\varepsilon\bigr]
 \le C_{\cK,\varepsilon}e^{-c_{\cK,\varepsilon}N}.
 \label{eq:tail}
\end{equation}
\end{corollary}

\begin{proof}
Let
\[
 Y_{N,s}
 \coloneqq\log\braket{e^{\rho_0NQ_{12}^2/2}}_s.
\]
By \eqref{eq:normalizedcoupling}, \eqref{eq:preconcentration} and \eqref{eq:lemreconcentration2},
\begin{equation}\label{eq:EY}
 \E Y_{N,s}=2NF_{N,s}(\rho_0)
 \le O_{\cK}(1)N\epsilon_N=o(N)
\end{equation}
uniformly in $s\in [0,1]$.
It is known that the map
\begin{align*}
((g_{ij})_{i,j}, (Z_i)_{i}) \mapsto Y_{N,s} = & \log
 \sum_{\sigma^1,\sigma^2}
 \exp\left(
 H_{N,s}(\sigma^1)+H_{N,s}(\sigma^2)
 +\frac{\rho _0 N}{2}Q_{12}^2\right) \\
&- \log
 \sum_{\sigma^1,\sigma^2}
 \exp\left(
 H_{N,s}(\sigma^1)+H_{N,s}(\sigma^2)
 \right) ,
\end{align*}
where we note that $H_{N,s}$ depends on $((g_{ij})_{i,j}, (Z_i)_{i})$, is Lipschitz continuous with Lipschitz constant $C_{\cK}\sqrt{N}$
(see \cite[Theorem~1.3.5]{Talagrand} and the discussion after the theorem).
The Gaussian concentration inequality
(see \cite[Theorem~1.3.4]{Talagrand})
%(see, for example,\cite{Ledoux})
therefore gives, for all
large enough $N$,
\[
 \mathbb P\left(
 |Y_{N,s} - \mathbb E Y_{N,s}|> t \right)
\le \exp \left( - N^{-1 } t^2/C_{\cK}  \right) \qquad (t\in [0,\infty )).
\]
This inequality and \eqref{eq:EY} yield
\begin{equation}\label{eq:tailY}
\mathbb P\left(
 Y_{N,s}>\frac{\rho_0\varepsilon^2N}{4}\right) \le C_{\cK,\varepsilon}e^{-c_{\cK,\varepsilon}N}.
\end{equation}
For each realization of the disorder, on $\{|Q_{12}|\ge\varepsilon\}$ we have
\(
\1_{\{|Q_{12}|\ge\varepsilon\}}
\le
\exp\left(
\frac{\rho_0N}{2}(Q_{12}^2-\varepsilon^2)
\right).
\) 
Therefore,
\[
\left\langle \1_{\{|Q_{12}|\ge\varepsilon\}}\right\rangle_s
\le
\exp\left(-\frac{\rho_0\varepsilon^2N}{2}\right)
\left\langle
\exp\left(\frac{\rho_0NQ_{12}^2}{2}\right)
\right\rangle_s
=
\exp\left(
-\frac{\rho_0\varepsilon^2N}{2}+Y_{N,s}
\right).
\]
Hence, 
\eqref{eq:tailY} implies
\begin{align*}
\nu_s\bigl[|Q_{12}|\ge\varepsilon\bigr]
&\le
\E\left[
\exp\left(
-\frac{\rho_0\varepsilon^2N}{2}+Y_{N,s}
\right)
\1_{\left\{Y_{N,s}\le \frac{\rho_0\varepsilon^2N}{4}\right\}}
\right]
+
\mathbb P\left(
Y_{N,s}>\frac{\rho_0\varepsilon^2N}{4}
\right)
\\
&\le
\exp\left(-\frac{\rho_0\varepsilon^2N}{4}\right)
+
\mathbb P\left(
Y_{N,s}>\frac{\rho_0\varepsilon^2N}{4}
\right).
\end{align*}
Therefore, we obtain \eqref{eq:tail}.
\end{proof}

\subsection{Cavity closure}\label{sec:cavityclosure}

Recall from \eqref{eq:ABC} that
\[
 A_s=\nu_s[Q_{12}^2],\qquad
 B_s=\nu_s[Q_{12}Q_{13}],\qquad
 C_s=\nu_s[Q_{12}Q_{34}].
\]
Introduce the three cavity modes
\begin{equation}
 U_s\coloneqq A_s-4B_s+3C_s,\qquad
 V_s\coloneqq2B_s-3C_s,\qquad
 D_s\coloneqq A_s-2B_s+C_s,
 \label{eq:cavity-modes}
\end{equation}
and put
\begin{equation}
 \kappa \coloneqq1-4q+3r,
 \qquad
 \zeta\coloneqq2q+q^2-3r.
 \label{eq:cavity-parameters}
\end{equation}

\begin{proposition}
\label{prop:cavity}
There are remainders
$\mathcal R^U_{N,s},\mathcal R^V_{N,s},\mathcal R^D_{N,s}$ such that
\begin{align}
 (1-s\beta^2\kappa )U_s
 &=\frac{\kappa }N+\mathcal R^U_{N,s},
 \label{eq:cavity-U}\\
 (1-s\beta^2\kappa )V_s-s\beta^2\zeta U_s
 &=\frac{\zeta}N+\mathcal R^V_{N,s},
 \label{eq:cavity-V}\\
 (1-s\alpha)D_s
 &=\frac{\alpha}{\beta^2N}+\mathcal R^D_{N,s},
 \label{eq:cavity-D}
\end{align}
and uniformly on $\cK\times[0,1]$,
\begin{equation}
 \max_{X\in\{U,V,D\}}|\mathcal R^X_{N,s}|
 \le C_{\cK}(N^{-3/2}+\nu_s[|Q_{12}|^3]).
 \label{eq:cavity-remainders}
\end{equation}
\end{proposition}

The proof follows from a straightforward, though technical, application of Talagrand's cavity argument; see \cite[Chapter~1]{Talagrand}. We postpone the details to Appendix~\ref{app:cavity-identities}.

Since $r\le q$,
\[
 \alpha-\beta^2\kappa
 =
 2\beta^2(q-r)\ge0.
\]
Thus $\beta^2\kappa \le\alpha<1$.  If $\kappa <0$, then
$1-s\beta^2\kappa \ge1$; if $\kappa \ge0$, then
$1-s\beta^2\kappa \ge1-\alpha$.  Hence \eqref{eq:delta} gives
\begin{equation}
 1-s\alpha\ge \delta _{\cK},
 \qquad
 1-s\beta^2\kappa \ge \delta _{\cK}
 \label{eq:denominator-bounds}
\end{equation}
for all $(\beta,h,s)\in\cK\times[0,1]$.
Equations \eqref{eq:cavity-U}--\eqref{eq:cavity-D}, in this order,
and the boundedness of $\zeta$ on $\cK$ imply
\[
 |U_s|+|V_s|+|D_s|
 \le C_{\cK}\left(\frac1N+N^{-3/2}+\nu_s[|Q_{12}|^3]\right).
\]
Since
\[
 A_s=-V_s-2U_s+3D_s,
\]
we conclude that
\begin{equation}
 A_s\le \frac{C_{\cK}}N
 +C_{\cK}\nu_s[|Q_{12}|^3].
 \label{eq:preabsorb}
\end{equation}

\section{Conclusions}\label{sec:conclusion}
\subsection{Absorption of the cubic remainder}\label{subsec:absorb}

Fix $\varepsilon>0$.  Since $|Q_{12}|\le2$,
\begin{align}
 \nu_s[|Q_{12}|^3]
 &\le
 \varepsilon\,\nu_s[Q_{12}^2]
 +8\,\nu_s[|Q_{12}|>\varepsilon]\nonumber\\
 &\le
 \varepsilon A_s
 +8C_{\cK,\varepsilon}e^{-c_{\cK,\varepsilon}N},
 \label{eq:thirdsplit}
\end{align}
where Corollary~\ref{cor:tail} was used.  Choose
$\varepsilon=\varepsilon_{\cK}$
so that the coefficient of $A_s$ obtained by substituting
\eqref{eq:thirdsplit} into \eqref{eq:preabsorb} is at most $1/2$.
We obtain
\[
 A_s=\nu_s[Q_{12}^2] \le\frac{C_{\cK}}N
\]
uniformly on $\cK\times[0,1]$, after increasing $C_{\cK}$ to include
the finitely many small values of $N$.  This proves the first claim
of Theorem~\ref{thm:main}.

\subsection{Free energy}\label{subsec:freeenergy}

Recall that $D_N(s):= P_s^* -\phi_N(s)$. The identity \eqref{eq:DNprime}, $D_N'(s) = \frac{\beta^2}{4}\nu_s[Q_{12}^2]$, with $P_1^*  =  \phi^{\mathrm{RS}}(\beta,h)$ and  $D_N(0)=0$ gives
\begin{equation}
 \phi^{\mathrm{RS}}(\beta,h)-\phi_N(1)
 =D_N(1)=
 \frac{\beta^2}{4}\int_0^1\nu_s[Q_{12}^2]\,\dd s.
 \label{eq:freeenergyidentity}
\end{equation}
The integrand is nonnegative, and the  
bound above makes the right-hand side at most $C_{\cK}/N$.  This proves the
free-energy claim of Theorem~\ref{thm:main}.

\subsection{Replicon susceptibility}\label{subsec:susceptibility}

The $O(N^{-1})$ second-moment estimate and
\eqref{eq:thirdsplit} imply, for every fixed $\varepsilon>0$,
\[
 \sup_{\cK\times[0,1]}
 N\nu_s[|Q_{12}|^3]
 \le C_{\cK}\varepsilon
 +8N C_{\cK,\varepsilon}e^{-c_{\cK,\varepsilon}N}.
\]
First let $N\to\infty$ and then $\varepsilon\downarrow0$.  Thus
\[
 \nu_s[|Q_{12}|^3]=o_{\cK}(N^{-1})
\]
uniformly in $s$.  Equation \eqref{eq:cavity-D}, together with
\eqref{eq:cavity-remainders} and \eqref{eq:denominator-bounds}, now gives
\[
 ND_s
 =
 \frac{\alpha}{\beta^2(1-s\alpha)}+o_{\cK}(1),
\]
which is the final claim of Theorem~\ref{thm:main}.

\appendix
\section{Proof of the cavity approximations}
\label{app:cavity-identities}

\begin{proof}[Proof of Proposition~\ref{prop:cavity}]
 All implicit constants in this proof depend only on $\cK$.  Write
\[
\eta_{N,s}\coloneqq N^{-3/2}+\nu_s[|Q_{12}|^3] ,\quad  \varepsilon_a\coloneqq\sigma_N^a,
 \quad
 Q^-_{ab}\coloneqq\frac1N\sum_{i<N}\sigma_i^a\sigma_i^b-q,
\]
so that
\begin{equation}
 Q_{ab}=Q^-_{ab}+\frac1N\varepsilon_a\varepsilon_b.
 \label{eq:Qab}
\end{equation}
By site exchangeability after averaging over the disorder,
\begin{align*}
 A_s&=\nu_s\!\left[Q_{12}(\varepsilon_1\varepsilon_2-q)\right],&
 B_s&=\nu_s\!\left[Q_{12}(\varepsilon_1\varepsilon_3-q)\right],&
 C_s&=\nu_s\!\left[Q_{12}(\varepsilon_3\varepsilon_4-q)\right].
\end{align*}

Use the last-spin cavity interpolation.  Write
$\sigma=(\rho,\varepsilon)$, with $\rho\in\{-1,1\}^{N-1}$, and replace
the interaction of the last spin with the first $N-1$ spins by
\[
 \frac{\beta\sqrt{su}}{\sqrt N}
 \sum_{i<N}\frac{g_{iN}+g_{Ni}}{\sqrt2}\rho_i\varepsilon
 +\beta\sqrt{s(1-u)q}\,\widehat z\,\varepsilon,
 \qquad 0\le u\le1,
\]
where $\widehat z$ is an independent standard Gaussian.  Denote the
corresponding averaged Gibbs expectation by $\nu_{s,u}$.  Then
$\nu_{s,1}=\nu_s$.  At $u=0$, the last spin is independent of the first
$N-1$ spins and is subject to the field
\[
 h+\beta\sqrt{(1-s)q}\,z_N+\beta\sqrt{sq}\,\widehat z
 \stackrel{\mathrm d}=h+\beta\sqrt q\,Z.
\]
Thus, for distinct replica indices,
\begin{equation}
 \mathbb E_0[\varepsilon_a\varepsilon_b]=q,
 \qquad
 \mathbb E_0[\varepsilon_a\varepsilon_b\varepsilon_c\varepsilon_d]=r.
 \label{eq:last-spin-moments-proof}
\end{equation}

Gaussian integration by parts gives, for every bounded function $F$ of
$n$ replicas,
\begin{align}
 \frac{\mathrm d}{\mathrm du}\nu_{s,u}[F]
 =s\beta^2\Bigg\{&
 \sum_{1\le a<b\le n}
 \nu_{s,u}[F\varepsilon_a\varepsilon_bQ^-_{ab}]
 -n\sum_{a=1}^n
 \nu_{s,u}[F\varepsilon_a\varepsilon_{n+1}Q^-_{a,n+1}]
 \nonumber\\
 &+\frac{n(n+1)}2
 \nu_{s,u}[F\varepsilon_{n+1}\varepsilon_{n+2}
 Q^-_{n+1,n+2}]
 \Bigg\}.
 \label{eq:cavity-derivative-proof}
\end{align}
Indeed, the derivative of the covariance of the interpolating
last-spin field is
$s\beta^2\varepsilon\varepsilon'Q^-(\rho,\rho')$, while its diagonal
value is independent of the configuration.

Set $M_3(u)\coloneqq\nu_{s,u}[|Q^-_{12}|^3]$.  Formula
\eqref{eq:cavity-derivative-proof} and $|Q^-_{ab}|\le2$ give
\[
 |M_3'(u)|\le C_{\cK}M_3(u).
\]
Since
\[
 |Q^-_{12}|^3\le4|Q_{12}|^3+\frac4{N^3},
\]
Gronwall's inequality, applied backward from $u=1$, yields
\begin{equation}
 \sup_{0\le u\le1}\nu_{s,u}[|Q^-_{12}|^3]
 \le C_{\cK}\eta_{N,s}.
 \label{eq:M3-gronwall-proof}
\end{equation}

For
\[
 P^-=Q^-_{ab}Q^-_{cd},
 \qquad (ab,cd)\in\{(12,12),(12,13),(12,34)\},
\]
the derivative of $\nu_{s,u}[P^-]$ is a finite linear combination of
expectations of three cavity overlaps.  Hence H\"older's inequality and
\eqref{eq:M3-gronwall-proof} show that
\[
 |\nu_{s,0}[P^-]-\nu_{s,1}[P^-]|
 \le C_{\cK}\eta_{N,s}.
\]
Moreover, \eqref{eq:Qab} gives
\[
 |P^--Q_{ab}Q_{cd}|
 \le\frac{|Q_{ab}|+|Q_{cd}|}{N}+\frac1{N^2}.
\]
Using
\[
 \frac{|x|}{N}\le\frac{|x|^3}{3}+\frac{2}{3N^{3/2}},
\]
we obtain
\begin{equation}
 \left|\nu_{s,0}[P^-]-\nu_s[Q_{ab}Q_{cd}]\right|
 \le C_{\cK}\eta_{N,s}.
 \label{eq:endpoint-replacement-proof}
\end{equation}

Define
\begin{align*}
 A_s^-&\coloneqq\nu_{s,0}[(Q^-_{12})^2],&
 B_s^-&\coloneqq\nu_{s,0}[Q^-_{12}Q^-_{13}],&
 C_s^-&\coloneqq\nu_{s,0}[Q^-_{12}Q^-_{34}],
\end{align*}
and form $U_s^-,V_s^-,D_s^-$ from $A_s^-,B_s^-,C_s^-$ as in
\eqref{eq:cavity-modes}.  It follows from
\eqref{eq:endpoint-replacement-proof} that
\begin{equation}
 |U_s^--U_s|+|V_s^--V_s|+|D_s^--D_s|
 \le C_{\cK}\eta_{N,s}.
 \label{eq:cavity-mode-replacement}
\end{equation}

Next write
\begin{align}
 A_s&=\nu_s[G_A]+\frac1N\nu_s[E_A],\nonumber\\
 B_s&=\nu_s[G_B]+\frac1N\nu_s[E_B],\nonumber\\
 C_s&=\nu_s[G_C]+\frac1N\nu_s[E_C],
 \label{eq:cavity-decomposition-proof}
\end{align}
where
\begin{align*}
 G_A&\coloneqq Q^-_{12}(\varepsilon_1\varepsilon_2-q),&
 E_A&\coloneqq\varepsilon_1\varepsilon_2
 (\varepsilon_1\varepsilon_2-q),\\
 G_B&\coloneqq Q^-_{12}(\varepsilon_1\varepsilon_3-q),&
 E_B&\coloneqq\varepsilon_1\varepsilon_2
 (\varepsilon_1\varepsilon_3-q),\\
 G_C&\coloneqq Q^-_{12}(\varepsilon_3\varepsilon_4-q),&
 E_C&\coloneqq\varepsilon_1\varepsilon_2
 (\varepsilon_3\varepsilon_4-q).
\end{align*}
Put
\[
 g_j(u)\coloneqq\nu_{s,u}[G_j],
 \qquad
 e_j(u)\coloneqq\nu_{s,u}[E_j],
 \qquad j\in\{A,B,C\},
\]
and take the same linear combinations as in \eqref{eq:cavity-modes}:
\begin{align*}
 g_U&\coloneqq g_A-4g_B+3g_C,&
 g_V&\coloneqq2g_B-3g_C,&
 g_D&\coloneqq g_A-2g_B+g_C,
\end{align*}
with analogous definitions for $e_U,e_V,e_D$.

Independence at $u=0$ gives $g_j(0)=0$.  Applying
\eqref{eq:cavity-derivative-proof} twice, H\"older's inequality and
\eqref{eq:M3-gronwall-proof} give
\begin{equation}
 g_X(1)=g_X'(0)+C_{\cK}\eta_{N,s},
 \qquad X\in\{U,V,D\}.
 \label{eq:cavity-Taylor-modes}
\end{equation}
For a replica edge $\{a,b\}$, \eqref{eq:last-spin-moments-proof}
implies
\begin{equation}
 \mathbb E_0[(\varepsilon_p\varepsilon_q-q)
 \varepsilon_a\varepsilon_b]
 =
 \begin{cases}
 1-q^2,&\{a,b\}=\{p,q\},\\
 q-q^2,&\{a,b\}\text{ shares one endpoint with }\{p,q\},\\
 r-q^2,&\{a,b\}\cap\{p,q\}=\varnothing.
 \end{cases}
 \label{eq:edge-rule-proof}
\end{equation}
Substitution into \eqref{eq:cavity-derivative-proof}, followed directly
by the three linear combinations above, gives
\begin{align}
 g_U'(0)&=s\beta^2\kappa  U_s^-,\nonumber\\
 g_V'(0)&=s\beta^2(\zeta U_s^-+\kappa V_s^-),
 \label{eq:scalar-first-derivatives}\\
 g_D'(0)&=s\alpha D_s^- .\nonumber
\end{align}
This is the cancellation for which the modes in
\eqref{eq:cavity-modes} were chosen.

The last-spin moments also give
\begin{align}
 e_U(0)&=(1-q^2)-4(q-q^2)+3(r-q^2)=\kappa ,\nonumber\\
 e_V(0)&=2(q-q^2)-3(r-q^2)=\zeta,
 \label{eq:scalar-diagonal-terms}\\
 e_D(0)&=1-2q+r=\frac{\alpha}{\beta^2}.
 \nonumber
\end{align}
Since $E_A,E_B,E_C$ are uniformly bounded,
\eqref{eq:cavity-derivative-proof} yields
\[
 |e_X'(u)|\le C_{\cK}\nu_{s,u}[|Q^-_{12}|],
 \qquad X\in\{U,V,D\}.
\]
The inequality
\[
 \frac{|Q^-_{12}|}{N}
 \le\frac{|Q^-_{12}|^3}{3}+\frac{2}{3N^{3/2}}
\]
and \eqref{eq:M3-gronwall-proof} therefore imply
\begin{equation}
 \frac1N|e_X(1)-e_X(0)|
 \le C_{\cK}\eta_{N,s},
 \qquad X\in\{U,V,D\}.
 \label{eq:scalar-diagonal-remainders}
\end{equation}

Taking the three mode combinations in
\eqref{eq:cavity-decomposition-proof}, and then using
\eqref{eq:cavity-mode-replacement},
\eqref{eq:cavity-Taylor-modes}, \eqref{eq:scalar-first-derivatives},
\eqref{eq:scalar-diagonal-terms}, and
\eqref{eq:scalar-diagonal-remainders}, gives
\begin{align*}
 U_s&=s\beta^2\kappa U_s+\frac{\kappa }N
 +C_{\cK}\eta_{N,s},\\
 V_s&=s\beta^2(\zeta U_s+\kappa V_s)+\frac{\zeta}N
 +C_{\cK}\eta_{N,s},\\
 D_s&=s\alpha D_s+\frac{\alpha}{\beta^2N}
 +C_{\cK}\eta_{N,s}.
\end{align*}
These are \eqref{eq:cavity-U}--\eqref{eq:cavity-remainders}.
\end{proof}

\section{Finite-step proof of the specialized Guerra--Talagrand bound}
\label{app:specialGT}



\subsection{A finite-step differentiation formula}

The following lemma isolates the only hierarchical Gaussian calculation
needed in the proof.

\begin{lemma}
\label{lem:finite-step-interpolation-derivative}
Let $\mathcal S$ be finite, let $\mu$ be a nonzero finite measure on
$\mathcal S$, and fix
\[
0=m_0\leq m_1\leq\cdots\leq m_M\leq m_{M+1}=1.
\]
Let $\{G_{k,t}(\xi):\xi\in\mathcal S,\ t\in[0,1]\}_{0\le k\le M}$ be independent centered Gaussian processes with almost surely $C^1$ sample paths in $t \in (0,1)$, and let $K_{k,t}$ denote their covariance kernels. Assume that $K_{k,t}$ is bounded.  Put
\[
\mathcal G_t(\xi) \coloneqq \sum_{k=0}^M G_{k,t}(\xi),
\qquad
C_t^{(j)}(\xi,\eta) \coloneqq \sum_{k=0}^jK_{k,t}(\xi,\eta),
\]
and write $C_t=C_t^{(M)}$. For a deterministic function
$V:\mathcal S\to\mathbb R$, define
\[
X_M(t)=\log\sum_{\xi\in \mathcal S}
\mu(\xi )e^{\mathcal G_t(\xi)+V(\xi)}
\]
and, recursively for $j=M,\ldots,1$,
\[
X_{j-1}(t) \coloneqq
\begin{cases}
\displaystyle
\frac1{m_j}\log\E_j e^{m_jX_j(t)},&m_j>0,\\[3mm]
\E_jX_j(t),&m_j=0,
\end{cases}
\]
where $\E_j$ integrates only the field $G_{j,t}$.  For
$0\leq j\leq M$, set
\[
\mathcal I_jA \coloneqq \frac{\E_j[Ae^{m_jX_j(t)}]}{\E_je^{m_jX_j(t)}},
\]
and let
\[
p_t(\xi) \coloneqq
\frac{\mu(\xi)e^{\mathcal G_t(\xi)+V(\xi)}}
{\sum_{\eta\in\mathcal S}
\mu(\eta)e^{\mathcal G_t(\eta)+V(\eta)}}.
\]
For deterministic functions $F$ and $0\leq j\leq M$, define
\begin{align*}
\E\langle F(\xi) \rangle_t
&\coloneqq \sum_{\xi\in\mathcal S}F(\xi)
\mathcal I_0\cdots\mathcal I_Mp_t(\xi),\\
\E\langle F(\xi,\eta)\rangle_{j,t}
&\coloneqq \sum_{\xi,\eta\in\mathcal S}F(\xi,\eta)
\mathcal I_0\cdots\mathcal I_j\left[
\bigl(\mathcal I_{j+1}\cdots\mathcal I_Mp_t(\xi)\bigr)
\bigl(\mathcal I_{j+1}\cdots\mathcal I_Mp_t(\eta)\bigr)
\right],
\end{align*}
where an empty operator product is the identity.  Then, for $0<t<1$,
\begin{equation}
\begin{aligned}
\frac{\mathrm d}{\mathrm dt} \E_0X_0(t)
={}&\frac1{2}\E\left\langle
\dot C_t(\xi,\xi)
\right\rangle_t-\frac1{2}\sum_{j=0}^M(m_{j+1}-m_j)
\E\left\langle
\dot C_t^{(j)}(\xi,\eta)
\right\rangle_{j,t}
\end{aligned}
\label{eq:finite-step-interpolation-derivative}
\end{equation}
where $\dot C_t^{(j)}(\xi,\eta) \coloneqq \frac{\partial}{\partial t} C_t^{(j)}(\xi,\eta)$.
\end{lemma}


\begin{proof}
Throughout the proof, for fixed $t$, we regard $X_j(t)$ as a smooth function of the Gaussian coordinates
\[
\{G_{k,t}(\xi):0\leq k\leq j,\ \xi\in\mathcal S\}
\]
that have not yet been integrated out; the fields at levels
$j+1,\ldots,M$ have already been averaged out by the recursion.
Thus $\partial_xX_{j-1}=\mathcal I_j(\partial_xX_j)$ whenever $x$ is
a Gaussian coordinate not integrated by $\E_j$ or the time coordinate $t$.  Indeed, if
$Z_j=\E_je^{m_jX_j}$ and $m_j>0$, then
\begin{align}
\partial_xX_{j-1}
&=\frac1{m_j}\frac{\partial_xZ_j}{Z_j}
=\frac{\E_j[e^{m_jX_j}\partial_xX_j]}{Z_j}
=\mathcal I_j(\partial_xX_j).
\label{eq:finite-step-first-chain-rule}
\end{align}
The same identity for $m_j=0$ follows by differentiating
$X_{j-1}=\E_jX_j$.
From \eqref{eq:finite-step-first-chain-rule} if follows that
\begin{equation}\label{eq:dX=IP}
\partial_{G_{k,t}(\xi)}X_k =\mathcal I_{k+1}\cdots\mathcal I_Mp_t(\xi).
\end{equation}

For any differentiable $A$, the quotient rule gives the more general
identity
\begin{align}
\partial_y\mathcal I_jA
={}&\mathcal I_j(\partial_yA)
+m_j\mathcal I_j(A\partial_yX_j)
-m_j\mathcal I_j(A)\mathcal I_j(\partial_yX_j).
\label{eq:derivative-tilted-expectation}
\end{align}
Consequently,
\begin{align}
\partial_{xy}X_{j-1}
={}&\mathcal I_j(\partial_{xy}X_j)
+m_j\left\{
\mathcal I_j(\partial_xX_j\partial_yX_j)
-\mathcal I_j(\partial_xX_j)
 \mathcal I_j(\partial_yX_j)
\right\}.
\label{eq:finite-step-second-chain-rule}
\end{align}
Again, this includes $m_j=0$, since the last two terms then vanish.

By \eqref{eq:finite-step-first-chain-rule}, we have
\begin{equation}
\begin{aligned}
    \frac{\mathrm d}{\mathrm dt}\E X_0 
&=\mathcal{I}_0\dots \mathcal{I}_{M}\frac{\mathrm d}{\mathrm dt} X_{M}(t)\\
&=\sum_{k=0}^M\mathcal{I}_0\dots \mathcal{I}_{M}\left( \sum_{\xi\in \mathcal{S}} p_t(\xi) \dot{G}_{k,t}(\xi)\right)\\
&= \sum_{k=0}^M \sum_{\xi\in \mathcal{S}} \mathcal{I}_0\dots\mathcal{I}_k\left(\dot{G}_{k,t}(\xi) (\mathcal{I}_{k+1} \dots\mathcal{I}_{M} (p_t(\xi)) )\right).
\label{eq:gaussian-covariance-differentiation}
\end{aligned}
\end{equation}
When $m_k>0$, by Gaussian integration by parts (see \cite[Vol. 1, Sec. A.4]{Talagrand}), i.e. for centered jointly  Gaussians $X,Y_i$,  
$$\mathbb E[Xf(Y_1\dots Y_d)] = 
\sum_{i=1}^d \operatorname{Cov}(X,Y_i),\mathbb E[\partial_i f(Y)],$$ 
where $\operatorname{Cov}(\dot{G}_{k,t}(\xi),\dot{G}_{k,t}(\eta)) = \dot K_{k,t}(\xi,\eta) $, by \eqref{eq:dX=IP} we have
\begin{align*}
 \mathcal{I}_k\left(\dot{G}_{k,t}(\xi) (\mathcal{I}_{k+1} \dots\mathcal{I}_{M} p_t(\xi) )\right)& = Z_k^{-1}  \mathbb E_k\left[e^{m_k X_k} \dot{G}_{k,t}(\xi) (\mathcal{I}_{k+1} \dots\mathcal{I}_{M} p_t(\xi) )\right] \\
 & = \sum_{\eta\in \mathcal{S}} \dot K_{k,t}(\xi,\eta) (Z_k)^{-1}\mathbb E_k\left[\partial_{{G}_{k,t}(\eta)}  ({e^{m_k X_k}} (\mathcal{I}_{k+1} \dots\mathcal{I}_{M} p_t(\xi) )) \right]\\
 &=\sum_{\eta\in \mathcal{S}} \dot K_{k,t}(\xi,\eta) (m_k Z_k)^{-1}\mathbb E_k\left[\partial_{{G}_{k,t}(\eta)} \partial_{{G}_{k,t}(\xi)} {e^{m_k X_k}}\right] .
\end{align*}
Let $H_{j,t}(\xi)=\sum^j_{r=0} G_{r,t}(\xi)$.   For fixed $k$, write
\[
u_{k}^\xi=\partial_{H_{k,t}(\xi)}X_k,
\qquad
h_{k}^{\xi\eta}
=\partial_{H_{k,t}(\xi)}\partial_{H_{k,t}(\eta)}X_k.
\]
The contribution obtained by differentiating only the
covariance of $G_{k,t}$ is
\begin{align*}
(m_k Z_k)^{-1}\E_k\!\left[
\partial_{G_k(\xi)}\partial_{G_k(\eta)}e^{m_kX_k}
\right]
&= 
\mathcal I_k\left(
h_k^{\xi\eta}+m_ku_k^\xi u_k^\eta
\right).
\end{align*}
For $m_k=0$, the same formula follows directly from
\eqref{eq:gaussian-covariance-differentiation}, with the second summand
equal to zero.  At $k=0$ it also holds because
$\mathcal I_0=\E_0$ and $m_0=0$. 
Hence, letting
\[
\mathcal P_k=\mathcal I_0\mathcal I_1\cdots\mathcal I_k,
\qquad
Q_k^{\xi\eta}=h_k^{\xi\eta}+m_ku_k^\xi u_k^\eta,
\]
by \eqref{eq:gaussian-covariance-differentiation} we have
\begin{equation}
\frac{\mathrm d}{\mathrm dt}\E X_0
=\frac12\sum_{k=0}^M\sum_{\xi,\eta\in\mathcal S}
\dot K_{k,t}(\xi,\eta)\mathcal P_kQ_k^{\xi\eta}.
\label{eq:finite-step-level-contributions}
\end{equation}

It remains to compute $\mathcal P_kQ_k^{\xi\eta}$.  At the terminal
level, direct differentiation gives
\[
u_M^\xi=p_t(\xi),\qquad
h_M^{\xi\eta}
=\mathbf 1_{\{\xi=\eta\}}p_t(\xi)-p_t(\xi)p_t(\eta).
\]
Since $m_{M+1}=1$, it follows that
\begin{equation}\label{eq:QM}
Q_M^{\xi\eta}
=\mathbf 1_{\{\xi=\eta\}}p_t(\xi)
-(m_{M+1}-m_M)u_M^\xi u_M^\eta.
\end{equation}
Note that \eqref{eq:dX=IP} implies $\partial_{H_{j-1,t}(\xi)}X_j = \partial_{H_{j,t}(\xi)}X_j$.
For $j=M,\ldots,1$, equations
\eqref{eq:finite-step-first-chain-rule} and
\eqref{eq:finite-step-second-chain-rule} yield
\begin{align*}
u_{j-1}^\xi
&=\mathcal I_ju_j^\xi, \quad 
h_{j-1}^{\xi\eta}
=\mathcal I_jh_j^{\xi\eta}
+m_j\left(
\mathcal I_j(u_j^\xi u_j^\eta)
-u_{j-1}^\xi u_{j-1}^\eta
\right),
\end{align*}
and hence
\[
Q_{j-1}^{\xi\eta}=h_{j-1}^{\xi\eta}+m_{j-1}u_{j-1}^\xi u_{j-1}^\eta
=\mathcal I_jQ_j^{\xi\eta}
-(m_j-m_{j-1})u_{j-1}^\xi u_{j-1}^\eta.
\]
By applying $\mathcal P_{j-1}$ to both sides of this equality, we have
\[
\mathcal P_{j-1}Q_{j-1}^{\xi\eta}
=\mathcal P_jQ_j^{\xi\eta}
-(m_j-m_{j-1})
\mathcal P_{j-1}(u_{j-1}^\xi u_{j-1}^\eta).
\]
Iterating from $j=M$ down to $j=k+1$ and using \eqref{eq:QM}
gives
\begin{align}
\mathcal P_kQ_k^{\xi\eta}
={}&\mathcal P_M
\left(\mathbf 1_{\{\xi=\eta\}}p_t(\xi)\right)
-\sum_{j=k}^M(m_{j+1}-m_j)
\mathcal P_j(u_j^\xi u_j^\eta).
\label{eq:finite-step-telescoped-hessian}
\end{align}
Note that similarly to \eqref{eq:dX=IP} it holds that $u_j^\xi =\mathcal I_{j+1}\cdots\mathcal I_Mp_t(\xi)$.
Substituting \eqref{eq:finite-step-telescoped-hessian} into
\eqref{eq:finite-step-level-contributions} therefore give
\begin{align*}
   \frac{\mathrm d}{\mathrm dt}\E X_0 
={}&\frac12\E\left\langle
\sum_{k=0}^M\dot K_{k,t}(\xi,\xi)
\right\rangle_t-\frac12\sum_{k=0}^M\sum_{j=k}^M(m_{j+1}-m_j)
\E\left\langle
\dot K_{k,t}(\xi,\eta)
\right\rangle_{j,t}\\
={}&\frac12\E\left\langle
\dot C_t(\xi,\xi)
\right\rangle_t
-\frac12\sum_{j=0}^M(m_{j+1}-m_j)
\E\left\langle
\dot C_t^{(j)}(\xi,\eta)
\right\rangle_{j,t},
\end{align*}
where we have used $\sum_{k=0}^M\sum_{j=k}^M\dot K_{k,t}(\xi,\eta)=\sum_{j=0}^M\sum_{k=0}^j \dot K_{k,t}(\xi,\eta)=\sum_{j=0}^M \dot C_t^{(j)}(\xi,\eta)$.
\end{proof}

\subsection{Application to the two-replica path}
Fix $N\geq1$, $\beta>0$, $h\in\mathbb R$, $s\in[0,1]$, and $q\in(0,1)$.
We use the same notation as in Section \ref{subsec:GT}.
Let $H_{N,s}^{\mathrm{cen}}$ be a centered Gaussian field on
$\Sigma_N$ normalized by
\begin{equation}
\E H_{N,s}^{\mathrm{cen}}(\sigma)
H_{N,s}^{\mathrm{cen}}(\tau)
=N\xi_s(R(\sigma,\tau)).
\label{eq:appendix-exact-covariance}
\end{equation}
Set
\[
H_{N,s}(\sigma)=H_{N,s}^{\mathrm{cen}}(\sigma)
+h\sum_{i=1}^N\sigma_i
\]
and
\[
\mathcal R_N=\{-1,-1+2/N,\ldots,1\}.
\]
For $v\in\mathcal R_N$, let
\begin{equation}
Z_{N,s}^{(2)}(v)=
\sum_{\substack{\sigma^1,\sigma^2\in\Sigma_N\\
R(\sigma^1,\sigma^2)=v}}
e^{H_{N,s}(\sigma^1)+H_{N,s}(\sigma^2)}.
\label{eq:appendix-constrained-partition}
\end{equation}

%We use the finite data directly.  Put $r=|v|$, let $\iota=1$ for
%$v\ge0$ and $\iota=-1$ for $v<0$, and let
%\[
%0=a_0<a_1<\cdots<a_M=1
%\]
%be the increasing list of the distinct elements of $\{0,q,r,1\}$.
%Set
%\begin{align}
%Q_j^v&=
%\begin{pmatrix}
%a_j&\iota(a_j\wedge r)\\
%\iota(a_j\wedge r)&a_j
%\end{pmatrix},
%\\
%m_j&=
%\begin{cases}
%0,&a_j\le q,\\[1mm]
%\frac12,&q<a_j\le r,\\[1mm]
%1,&a_j>\max\{q,r\},
%\end{cases}
%\qquad 1\le j\le M,
%\label{eq:appendix-hierarchy-levels}\\
%B_j&=\beta^2(1-s)q
%\begin{pmatrix}1&1\\1&1\end{pmatrix}+s\beta^2Q_j^v,
%\qquad 0\le j\le M,
%\\
%d_j&=\frac{s\beta^2}{2}
%\sum_{k,\ell =1}^2(Q_{j,k\ell}^v)^2,
%\qquad 0\le j\le M,
%\end{align}
%where $Q_{j,k\ell}^v$ is the $(k,\ell)$ entry of $Q_{j}^v$.
%Here, note that if $q\geq r$, $m_j$ takes values $0$ or $1$.
%As in \eqref{eq:GT-overlap-increments}, the overlap increments are
%nonnegative multiples either of
%$\left(\begin{smallmatrix}1&\iota\\\iota&1\end{smallmatrix}\right)$
%or of $I_2$.  Hence
%\begin{equation}
%B_j-B_{j-1}=s\beta^2(Q_j^v-Q_{j-1}^v)\succeq0.
%\label{eq:finite-step-vector-increments}
%\end{equation}
%Moreover, $d_j=d(a_j)$, where
%\[
%d(a)=
%\begin{cases}
%2s\beta^2a^2,&0\le a\le r,\\
%s\beta^2(a^2+r^2),&r\le a\le1.
%\end{cases}
%\]
%Thus
%\begin{equation}
%d_j-d_{j-1}\ge0.
%\label{eq:finite-step-scalar-increments}
%\end{equation}
%The sequence in \eqref{eq:appendix-hierarchy-levels} is nondecreasing.
%This remains true when $r=q$ or $r=1$; in the latter case its largest
%value is $1/2$.
%
%For $\lambda\in\mathbb R$, define $U_s^{\lambda,v}$ from
%$f_\lambda$ in \eqref{eq:GTterminal} by the recursion
%\eqref{eq:U-stage-recursion} with the
%data above.  Let $Z$ be standard normal, put
%$Y_0=h+\beta\sqrt{(1-s)q}\,Z$, and set
%\[
%\mathcal L_s(v)=\frac12\sum_{j=1}^Mm_j(d_j-d_{j-1}),
%\qquad
%\mathfrak P_s(\lambda,v)=
%2\log2+\E U_s^{\lambda,v}(Y_0,Y_0)
%-\lambda v-\mathcal L_s(v).
%\]

We prove \eqref{eq:GTfunctionalbound} as follows.

\begin{lemma}
\label{lem:specialGT}
For every $N\geq1$ and $v\in\mathcal R_N$,
\begin{equation}
\frac1N\E\log Z_{N,s}^{(2)}(v)
\leq\inf_{\lambda\in\mathbb R}\mathfrak P_s(\lambda,v).
\label{eq:appendix-GT-infimum}
\end{equation}
\end{lemma}

\begin{proof}
%Fix $\lambda\in\mathbb R$.  The monotonicity of the levels required
%by Lemma~\ref{lem:finite-step-interpolation-derivative}, as well as the
%positivity of all Gaussian variances below, was verified in
%\eqref{eq:finite-step-vector-increments} and
%\eqref{eq:finite-step-scalar-increments}.
In view of \eqref{eq:finite-step-vector-increments} and \eqref{eq:finite-step-scalar-increments}, we are able to choose independent Gaussian vectors $z_{i,j}\in\mathbb R^2$ and
independent Gaussian variables $y_j$ with
\[
\E z_{i,j}z_{i,j}^{\mathsf T}=B_j-B_{j-1},
\qquad
\E y_j^2=d_j-d_{j-1}.
\]
Let $Y_{i,0}$ be independent copies of
$Y_0=h+\beta\sqrt{(1-s)q}\,Z$. Since
\[
B_0=\beta^2(1-s)q
\begin{pmatrix}1&1\\1&1\end{pmatrix},
\]
the vector $(Y_{i,0}-h,Y_{i,0}-h)$ supplies the level-zero covariance $B_0$.

Write
\[
{\mathcal{S}}_N(v)=\{(\sigma^1,\sigma^2)\in\Sigma_N^2:
R(\sigma^1,\sigma^2)=v\}.
\]
For $\boldsymbol\sigma=(\sigma^1,\sigma^2)\in{\mathcal{S}}_N(v)$, introduce
the independent Gaussian random variables $\{ G_{j,t}(\boldsymbol\sigma)\} _{j=0}^M$ with parameter $t \in [0,1]$ as
\begin{align*}
G_{0,t}(\boldsymbol\sigma)
\coloneqq {}&\sqrt t\bigl(H_{N,s}^{\mathrm{cen}}(\sigma^1)
+H_{N,s}^{\mathrm{cen}}(\sigma^2)\bigr)+\sqrt{1-t}\sum_{i=1}^N\sum_{a=1}^2
(Y_{i,0}-h)\sigma_i^a,\\
G_{j,t}(\boldsymbol\sigma)
\coloneqq {}&\sqrt{1-t}\sum_{i=1}^N\sum_{a=1}^2z_{i,j}^a\sigma_i^a
+\sqrt{tN}\,y_j,
\qquad 1\leq j\leq M,
\end{align*}
and the deterministic term
\[
V(\boldsymbol\sigma) \coloneqq h\sum_{i=1}^N(\sigma_i^1+\sigma_i^2)
+\lambda\sum_{i=1}^N\sigma_i^1\sigma_i^2-\lambda Nv.
\]
Let
\[
X_M(t) \coloneqq \log\sum_{\boldsymbol\sigma\in{\mathcal{S}}_N(v)}
\exp\left(\sum_{j=0}^MG_{j,t}(\boldsymbol\sigma)
+V(\boldsymbol\sigma)\right),
\]
define $\{ X_j(t)\} _{j=0}^M$ as in Lemma~\ref{lem:finite-step-interpolation-derivative},
and set
\[
\varphi_N(t) \coloneqq \frac1N\E_0X_0(t).
\]

We first compute the values at the endpoints.  At $t=0$, remove the overlap
constraint from the partition function.  The maps in
\eqref{eq:U-stage-recursion} are monotone, and the independent site
increments factorize.  The defining recursion {\color{red}What is this?} therefore gives
\begin{equation}
\varphi_N(0)
\leq2\log2+\E U_s^{\lambda,v}(Y_0,Y_0)-\lambda v.
\label{eq:GT-endpoint-zero}
\end{equation}
The factor $2\log2$ comes from the normalization $1/4$ in
$f_\lambda$.

At $t=1$, the Lagrange term vanishes on ${\mathcal{S}}_N(v)$, while the
sitewise Gaussian increments vanish. For a quantity $x$ independent
of $y_j$, the $j$th recursion gives
\[
\begin{cases}
\displaystyle
\frac1{m_j}\log\E_j e^{m_j(x+\sqrt N\,y_j)}
=x+\frac{Nm_j}{2}(d_j-d_{j-1}),&m_j>0,\\[3mm]
\displaystyle
\E_j(x+\sqrt N\,y_j)=x,&m_j=0.
\end{cases}
\]
Hence, by applying this equality recursively we have {\color{red} ($m_M=1$? Is the label should be shifted?)}
\begin{align}
\varphi_N(1)
&=\frac1N E_0 \dots E_{M-1} \left[ \log \E _M e^{m_M X_M(1)} \right] 
+\frac12\sum_{j=1}^Mm_j(d_j-d_{j-1})
\notag\\
&=\frac1N\E\log Z_{N,s}^{(2)}(v)+\mathcal L_s(v)
\label{eq:GT-endpoint-one}
\end{align}
where $Z_{N,s}^{(2)}(v)$  and $\mathcal L_s(v)$ are defined in \eqref{eq:appendix-constrained-partition} and \eqref{eq:defLv}, respectively.

It remains to compute the derivative of $\varphi _N$. Set $m_0=0$, $m_{M+1}=1$, and
for two spin pairs $\boldsymbol\sigma,\boldsymbol\tau\in{\mathcal{S}}_N(v)$
write
\[
q_j^{k\ell}=Q_{j,k\ell}^v,
\qquad
R^{k\ell}=R(\sigma^k,\tau^\ell).
\]
Also write
\[
\theta_s(x)=x\xi_s'(x)-\xi_s(x)=\frac{s\beta^2}{2}x^2.
\]
Denote the covariance of $\sum _{i=0}^j G_{i,t} (\boldsymbol\sigma)$ and $\sum _{i=0}^j G_{i,t} (\boldsymbol\tau)$ by $C_t^{(j)}(\boldsymbol\sigma,\boldsymbol\tau)$.
If $\boldsymbol\sigma _i = \boldsymbol\tau _i$ for $i=1,2,\dots ,j$, then it holds that
\begin{align}
C_t^{(j)}(\boldsymbol\sigma,\boldsymbol\tau)&
={}tN\sum_{k,\ell=1}^2\xi_s(R^{k\ell})
+(1-t)N\sum_{k,\ell=1}^2\xi_s'(q_j^{k\ell})R^{k\ell}
 +tN\sum_{k,\ell=1}^2\theta_s(q_j^{k\ell})
\label{eq:GT-split-covariance} {\color{red}Why?}\\
&=N\sum_{k,\ell=1}^2\xi_s'(q_j^{k\ell})R^{k\ell}+ t\left(\frac{s \beta^2 N}{2}\sum_{k,\ell=1}^2(R^{k\ell}-q_j^{k\ell})^2\right),
\end{align}
due to $\xi_s(x)-\xi_s(y)-\xi_s'(y)(x-y)=\frac{s\beta^2}{2}(x-y)^2$. 
For a single constrained pair, its $2\times2$ overlap matrix is
\[
\bigl(R(\sigma^k,\sigma^\ell )\bigr)_{k,\ell =1}^2
=\begin{pmatrix}1&v\\v&1\end{pmatrix}=Q_M^v.
\]
Consequently,
\[
\dot C_t^{(M)}(\boldsymbol\sigma,\boldsymbol\sigma)
=\frac{s\beta^2 N}{2}\sum_{k,\ell=1}^2
(Q_{M,k\ell}^v-Q_{M,k\ell}^v)^2=0.
\]
Hence, Lemma~\ref{lem:finite-step-interpolation-derivative} 
now yields
\begin{equation}
\varphi_N'(t)
=-\frac{s\beta^2}{4}
\sum_{j=0}^M(m_{j+1}-m_j)
\E\left\langle
\sum_{k,\ell=1}^2(R^{k\ell}-Q_{j,k\ell}^v)^2
\right\rangle_{j,t}
\leq0.
\label{eq:special-GT-interpolation-derivative}
\end{equation}
Combining \eqref{eq:GT-endpoint-zero},
\eqref{eq:GT-endpoint-one}, and
\eqref{eq:special-GT-interpolation-derivative}, and taking the infimum over $\lambda$ proves
\eqref{eq:appendix-GT-infimum}.
\end{proof}


\section{Proof of a technical lemma}\label{app: technical}
\begin{proof}[Proof of Lemma~\ref{lem:GTcontinuity}]
Fix $L<\infty$ and let
$(\beta_n,h_n,s_n,\lambda_n,v_n)\to(\beta,h,s,\lambda,v)$, with
$|\lambda_n|\le L$.  Write $q_n=q(\beta_n,h_n)$ and $r_n=|v_n|$.
By the continuity of $q$ established after \eqref{eq:q}, $q_n\to q$,
while $r_n\to r\coloneqq|v|$.  There are only finitely many possible
orders of $q_n$ and $r_n$ and signs of $v_n$, so it is enough to work
on a subsequence on which these choices are fixed.

For each $n$, the definition of $U_{s_n}^{\lambda_n,v_n}$ is a
composition of at most three operators from \eqref{eq:GTsemigroup}.
On a stage whose limiting breakpoint length is positive, its level is
eventually fixed and the covariance
$B_{j,n}-B_{j-1,n}$ converges to its limiting covariance.  A stage
whose breakpoints coalesce has covariance tending to zero, and its
operator converges to the identity even if its level changes.  This
accounts for $v=0$, $|v|=q$, and $|v|=1$.  If $s=0$, all the
covariances vanish, so the same conclusion is immediate.

On $|\lambda|\le L$, the terminal data obey
\[
 |f_\lambda(x_1,x_2)|\le |x_1|+|x_2|+L,
 \qquad |\partial_\lambda f_\lambda|\le1.
\]
For $m>0$, differentiation of \eqref{eq:GTsemigroup} gives
\[
 \partial_\lambda(\mathcal T_{m,C}F)
 =\frac{\E[(\partial_\lambda F)(x+G_C)e^{mF(x+G_C)}]}
        {\E e^{mF(x+G_C)}},
\]
with the ordinary Gaussian expectation when $m=0$.  Hence every step
preserves $|\partial_\lambda F|\le1$.  The analogous coordinate
derivative formula and $|\partial_{x_a}f_\lambda|\le1$ give a common
local equicontinuity bound.  For any convergent sequence of stage
covariances $C_n\to C$, couple
$G_{C_n}=C_n^{1/2}Z$.  Dominated convergence shows that each operator
and its $\lambda$-derivative converge locally uniformly.  The same
argument applies when $C=0$, independently of the value of $m_{j,n}$.
Iterating over the finite composition gives
local uniform convergence of both $U_{s_n}^{\lambda_n,v_n}(\cdot)$
and its $\lambda$-derivative.

Finally, as $n\to \infty$
\[
 Y_{0,n}=h_n+\beta_n\sqrt{(1-s_n)q_n}\,Z\longrightarrow
 Y_0=h+\beta\sqrt{(1-s)q}\,Z
 \quad\text{in }L^1,
\]
and the sequence $(Y_{0,n})$ is uniformly integrable.  The local
uniform convergence above and the common linear-growth bound therefore
give convergence after evaluation at $(Y_{0,n},Y_{0,n})$.  Since
\eqref{eq:GTcorrection} holds and $P_s^*$ is continuous, the
definition \eqref{eq:GTfunctional} and
\[
 \partial_\lambda\mathfrak P_s(\lambda,v)
 =\E\partial_\lambda U_s^{\lambda,v}(Y_0,Y_0)-v
\]
prove the asserted joint continuity.
\end{proof}
\section*{Acknowledgements}

S.K.\ acknowledges support from JSPS KAKENHI Grant Numbers JP22H00099 and JP23K20801.
S.N.\ acknowledges support from JSPS KAKENHI Grant Numbers JP24K16937 and JP25K00911.

\section*{Statement on the Use of Artificial Intelligence}

Codex was used to assist with revising the presentation and performing compilation checks. The authors take full responsibility for the content of this work.

\begin{thebibliography}{10}

\bibitem{AizenmanLebowitzRuelle1987}
M.~Aizenman, J.~L.~Lebowitz, and D.~Ruelle,
\newblock
\emph{Some rigorous results on the Sherrington--Kirkpatrick spin glass model},
\newblock \emph{Communications in Mathematical Physics} 112 (1987), 3--20.

\bibitem{AmitGutfreundSompolinsky1985}
D.~J.~Amit, H.~Gutfreund, and H.~Sompolinsky,
\newblock
\emph{Spin-glass models of neural networks},
\newblock \emph{Physical Review A} 32 (1985), 1007--1018.

\bibitem{deAlmeidaThouless1978}
J.~R.~L.~de Almeida and D.~J.~Thouless,
\newblock
\emph{Stability of the Sherrington--Kirkpatrick solution},
\newblock \emph{Journal of Physics A: Mathematical and General} 11 (1978), 983--990.

\bibitem{EdwardsAnderson1975}
S.~Edwards and P.~W.~Anderson,
\newblock
\emph{Theory of spin glasses},
\newblock \emph{Journal of Physics F: Metal Physics} 5 (1975), 965--974.

\bibitem{Hopfield1982}
J.~J.~Hopfield,
\newblock
\emph{Neural networks and physical systems with emergent collective computational abilities},
\newblock \emph{Proceedings of the National Academy of Sciences} 79 (1982), 2554--2558.

\bibitem{MezardMontanari2009}
M.~M{\'e}zard and A.~Montanari,
\newblock
\emph{Information, Physics, and Computation},
\newblock Oxford University Press, 2009.

\bibitem{MezardParisiZecchina2002}
M.~M{\'e}zard, G.~Parisi, and R.~Zecchina,
\newblock
\emph{Analytic and algorithmic solution of random satisfiability problems},
\newblock \emph{Science} 297 (2002), 812--815.

\bibitem{Nishimori2001}
H.~Nishimori,
\newblock
\emph{Statistical Physics of Spin Glasses and Information Processing: An Introduction},
\newblock Oxford University Press, 2001.

\bibitem{Parisi1979}
G.~Parisi,
\newblock
\emph{Infinite number of order parameters for spin glasses},
\newblock \emph{Physical Review Letters} 43 (1979), 1754--1756.

\bibitem{Parisi1980}
G.~Parisi,
\newblock
\emph{A sequence of approximated solutions to the S-K model for spin glasses},
\newblock \emph{Journal of Physics A: Mathematical and General} 13 (1980), L115--L121.

\bibitem{SherringtonKirkpatrick1975}
D.~Sherrington and S.~Kirkpatrick,
\newblock
\emph{Solvable model of a spin-glass},
\newblock \emph{Physical Review Letters} 35 (1975), 1792--1796.

\bibitem{ZdeborovaKrzakala2016}
L.~Zdeborov{\'a} and F.~Krzakala,
\newblock
\emph{Statistical physics of inference: Thresholds and algorithms},
\newblock \emph{Advances in Physics} 65 (2016), 453--552.

\bibitem{ChenGT}
W.-K.~Chen,
\newblock
\emph{Variational representations for the Parisi functional and the
two-dimensional Guerra--Talagrand bound},
\newblock \emph{Annals of Probability} 45 (2017), 3929--3966.
\href{https://arxiv.org/abs/1501.06635}{arXiv:1501.06635}.

\bibitem{DeyWu}
P.~S.~Dey and Q.~Wu,
\newblock
\emph{Fluctuation results for multi-species
Sherrington--Kirkpatrick model in the replica symmetric regime},
\newblock \emph{Journal of Statistical Physics} 185 (2021),
Paper No.~22.
\href{https://arxiv.org/abs/2012.13381}{arXiv:2012.13381}.

\bibitem{KusuokaNakajima1}
\newblock
S.~Kusuoka and S.~Nakajima,
\emph{A note on Lata{\l}a's argument in SK model},
\newblock preprint (2026).
\href{https://arxiv.org/abs/2607.23427}{arXiv:2607.23427}.

\bibitem{Ledoux}
M.~Ledoux,
\newblock
\emph{The Concentration of Measure Phenomenon},
\newblock Mathematical Surveys and Monographs, Vol.~89,
American Mathematical Society, Providence, RI, 2001.

\bibitem{Lopatto}
P.~Lopatto,
\newblock
\emph{Replica symmetry up to the de Almeida--Thouless line in the
Sherrington--Kirkpatrick model},
\newblock preprint (2026).
\href{https://arxiv.org/abs/2604.11921}{arXiv:2604.11921}.

\bibitem{Panchenko}
D.~Panchenko,
\newblock
\emph{The Sherrington--Kirkpatrick Model},
\newblock Springer Monographs in Mathematics, Springer, New York, 2013.

\bibitem{Talagrand}
M.~Talagrand,
\newblock
\emph{Mean Field Models for Spin Glasses}, Vols.~I--II,
\newblock Springer, 2011.

\end{thebibliography}


\end{document}

    
